\documentclass[10pt,letterpaper]{book}
\usepackage[margin=0.22in,bindingoffset=0.25in,includeheadfoot]{geometry}
\setlength{\headheight}{14pt}
\usepackage[T1]{fontenc}
\usepackage[utf8]{inputenc}
\usepackage{lmodern}
\usepackage{microtype}
\usepackage{xcolor}
\usepackage{listings}
\usepackage{booktabs}
\usepackage{emptypage}
\usepackage[colorlinks,linkcolor=refcol,urlcolor=refcol]{hyperref}
\definecolor{refcol}{RGB}{20,60,130}
\makeatletter                            % TOC: room for two-digit section numbers
\renewcommand*\l@section{\@dottedtocline{1}{1.5em}{3.2em}}
\makeatother

\definecolor{kw}{RGB}{20,60,130}
\definecolor{cm}{RGB}{80,120,80}
\definecolor{st}{RGB}{150,50,50}
\definecolor{rulecol}{RGB}{170,175,190}
\definecolor{numcol}{RGB}{150,150,150}

\lstdefinestyle{code}{language=Python,basicstyle=\ttfamily\small,
  keywordstyle=\color{kw}\bfseries,commentstyle=\color{cm}\itshape,
  stringstyle=\color{st},showstringspaces=false,keepspaces=true,
  columns=fullflexible,breaklines=true,breakatwhitespace=true,
  postbreak=\mbox{\textcolor{numcol}{$\hookrightarrow$}\space},
  upquote=true,aboveskip=6pt,belowskip=6pt,xleftmargin=0.5em,
  frame=leftline,framerule=0.8pt,rulecolor=\color{rulecol}}
\lstdefinestyle{ex}{basicstyle=\ttfamily\small,keepspaces=true,
  columns=fullflexible,breaklines=true,breakatwhitespace=true,
  postbreak=\mbox{\textcolor{numcol}{$\hookrightarrow$}\space},
  upquote=true,aboveskip=4pt,belowskip=7pt,xleftmargin=0.5em,
  frame=leftline,framerule=0.8pt,rulecolor=\color{rulecol},language={}}
\lstdefinestyle{file}{language=Python,basicstyle=\ttfamily\footnotesize,
  keywordstyle=\color{kw}\bfseries,commentstyle=\color{cm}\itshape,
  stringstyle=\color{st},showstringspaces=false,keepspaces=true,
  columns=fullflexible,breaklines=true,upquote=true,numbers=left,
  numberstyle=\tiny\color{numcol},numbersep=7pt,aboveskip=5pt,belowskip=5pt}

\setcounter{tocdepth}{1}
\setlength{\parskip}{2pt}
\title{\Huge\bfseries Haskell.py Teacher's Edition\\[6pt]\Large A Complete Course in Idiomatic Functional Python\\[14pt]\normalsize nine chapters, 41 homework problems, 10 compare-and-contrast apps}
\author{J.~L.~Clements~III}
\date{\today}
\begin{document}
\frontmatter
\maketitle

\chapter{How to Take This Course}
This book teaches \texttt{haskell.py} -- prelude.py's production sibling.
Where prelude.py reimplements every tool from scratch for pedagogy,
haskell.py WRAPS the standard library wherever it already has the tool
(\texttt{scanl} is \texttt{itertools.accumulate}, \texttt{memo} is
\texttt{functools.cache}) -- one vocabulary, zero reimplementation. Failure
is a value throughout: a true \texttt{NOTHING} sentinel keeps \texttt{None}
legal as everyday data, parsers fail with \texttt{FAIL}, and \texttt{Either}
carries WHY.

The whole library is one listing, up front, in Chapter 1 -- read it once,
top to bottom, the way you would read any short source file. The nine
chapters after it are not a second pass over the same material: each opens
with why its slice of the library exists and how its pieces fit together,
then goes straight to \textbf{Homework} -- drills cement each tool, applies
combine them, and challenges are interview-grade. Write each solution in a
scratch file with the given doctests pasted in; \texttt{python3 -m doctest}
is the referee.

\tableofcontents
\mainmatter
\part{The Nine Chapters}
\chapter{haskell.py, Complete}
Every name this book teaches, in one file, in the order it is defined --
docstrings and comments stripped to the section headers, which are the
only thing worth keeping as a wayfinding aid in a listing this short. Read
it once before the first chapter; the chapters that follow assume you have.
\begin{lstlisting}[style=file]
from functools import reduce as _reduce, partial, cache as memo, cmp_to_key
from itertools import accumulate as _acc, islice as _islice, pairwise as _pairwise
from itertools import chain, count, cycle, repeat
from operator import add, sub, mul, not_, eq
from math import gcd, lcm, hypot, isqrt, prod as product
from collections import Counter, defaultdict, deque
from heapq import heappush, heappop, merge as _hmerge
import re as _re
# 1. functions
identity = lambda x: x
const = lambda x: lambda *_: x
flip = lambda f: lambda a,b: f(b,a)
curry = lambda f: lambda x: lambda y: f(x,y)
uncurry = lambda f: lambda p: f(*p)
def compose(*fs):
    def g(x):
        for f in reversed(fs):
            x = f(x)
        return x
    return g
o = compose
def pipe(x,*fs):
    for f in fs:
        x = f(x)
    return x
on = lambda f,g: lambda x,y: f(g(x),g(y))
# 2. pairs
fst = lambda p: p[0]
snd = lambda p: p[1]
swap = o(tuple,reversed)
# 3. sentinels
class _Sentinel:
    __slots__ = ("_n",)
    def __init__(self,n): self._n = n
    def __repr__(self): return self._n
    def __bool__(self): return False
NOTHING = _Sentinel("NOTHING")
FAIL = _Sentinel("FAIL")
_MISS = _Sentinel("_MISS")
# 4. arithmetic & logic
div = lambda a,b: a // b
mod = lambda a,b: a % b
halve = lambda n: n // 2
even = lambda n: n % 2 == 0
odd = lambda n: n % 2 == 1
otherwise = True
quot = lambda a,b: -(-a // b) if (a < 0) != (b < 0) else a // b
rem = lambda a,b: a - b * quot(a,b)
quotRem = lambda a,b: (quot(a,b), rem(a,b))
signum = lambda x: (x > 0) - (x < 0)
succ = lambda x: chr(ord(x) + 1) if isinstance(x,str) else x + 1
pred = lambda x: chr(ord(x) - 1) if isinstance(x,str) else x - 1
# 5. folds, scans, unfolds
foldl = lambda f,xs,base=_MISS: _reduce(f,xs) if base is _MISS else _reduce(f,xs,base)
foldr = lambda f,xs,base: _reduce(flip(f),reversed(list(xs)),base)
foldl1 = foldl
foldr1 = lambda f,xs: _reduce(flip(f),reversed(list(xs)))
scanl = lambda f,xs,base=_MISS: list(_acc(xs,f) if base is _MISS else _acc(xs,f,initial=base))
scanl1 = scanl
scanr = lambda f,xs,base: scanl(flip(f),reversed(list(xs)),base)[::-1]
scanr1 = lambda f,xs: scanl(flip(f),reversed(list(xs)))[::-1]
def unfoldr(f,seed):
    out = []
    while (r := f(seed)) is not NOTHING:
        out.append(r[0]); seed = r[1]
    return out
def until(cond,f,x):
    while not cond(x):
        x = f(x)
    return x
replicateM = lambda n,xs: foldl(lambda acc,_: [s + [x] for s in acc for x in xs],range(n),[[]])
subsequences = lambda xs: foldl(lambda acc,x: acc + [s + [x] for s in acc],list(xs),[[]])
# 6. streams
def iterate(f,x):
    while True:
        yield x
        x = f(x)
take = lambda n,xs: list(_islice(iter(xs),n))
# 7. list basics
head = lambda xs: xs[0]
tail = lambda xs: xs[1:]
init = lambda xs: xs[:-1]
last = lambda xs: xs[-1]
drop = lambda n,xs: list(xs)[max(n,0):]
splitAt = lambda n,xs: (xs[:n], xs[n:]) if n >= 0 else (xs[:0], xs)
replicate = lambda n,x: [x] * n
elem = lambda x,xs: x in xs
notElem = lambda x,xs: x not in xs
null = lambda xs: len(xs) == 0
nub = o(list,dict.fromkeys)
enum = lambda xs,start=0: list(enumerate(xs,start))
pairwise = lambda xs: list(_pairwise(xs))
zip_ = lambda a,b: list(zip(a,b))
zip3 = lambda a,b,c: list(zip(a,b,c))
unzip = lambda ps: tuple(map(list,zip(*ps))) if ps else ([], [])
isPrefixOf = lambda p,xs: list(xs[:len(p)]) == list(p)
isSuffixOf = lambda p,xs: list(xs[len(xs) - len(p):]) == list(p)
lookup = lambda k,pairs: next((v for kk, v in pairs if kk == k),NOTHING)
stripPrefix = lambda p,xs: xs[len(p):] if isPrefixOf(p,xs) else NOTHING
# 8. higher-order lists
map_ = lambda f,xs: [f(x) for x in xs]
filter_ = lambda cond,xs: [x for x in xs if cond(x)]
concat = lambda xss: [x for xs in xss for x in xs]
concatMap = lambda f,xs: [y for x in xs for y in f(x)]
cross = lambda a,b: [(x, y) for x in a for y in b]
starmap = lambda f,pairs: [f(*p) for p in pairs]
zipWith = lambda f,a,b: [f(x,y) for x, y in zip(a,b)]
zipWith3 = lambda f,a,b,c: [f(x,y,z) for x, y, z in zip(a,b,c)]
find = lambda cond,xs: next((x for x in xs if cond(x)),NOTHING)
def partition(cond,xs):
    yes, no = [], []
    for x in xs:
        (yes if cond(x) else no).append(x)
    return (yes, no)
# 9. slicing & spans
def span(cond,xs):
    xs = list(xs)
    i = next((i for i, x in enumerate(xs) if not cond(x)),len(xs))
    return (xs[:i], xs[i:])
break_ = lambda cond,xs: span(o(not_,cond),xs)
def takewhile(cond,xs):
    for x in xs:
        if not cond(x): return
        yield x
def dropwhile(cond,xs):
    it = iter(xs)
    for x in it:
        if not cond(x):
            yield x
            break
    yield from it
def takeuntil(cond,xs):
    for x in xs:
        yield x
        if cond(x): return
takeWhile = lambda cond,xs: list(takewhile(cond,xs))
dropWhile = lambda cond,xs: list(dropwhile(cond,xs))
inits = lambda xs: [xs[:i] for i in range(len(xs) + 1)]
tails = lambda xs: [xs[i:] for i in range(len(xs) + 1)]
def chunksOf(n,xs):
    xs = list(xs)
    return [xs[i:i + n] for i in range(0,len(xs),n)]
def windows(n,xs):
    xs = list(xs)
    return [xs[i:i + n] for i in range(len(xs) - n + 1)]
def groupBy(eq,xs):
    out = []
    for x in xs:
        if out and eq(out[-1][0],x): out[-1].append(x)
        else: out.append([x])
    return out
group = partial(groupBy,eq)
def transpose(rows):
    rows = [list(r) for r in rows]
    out, i = [], 0
    while any(i < len(r) for r in rows):
        out.append([r[i] for r in rows if i < len(r)])
        i += 1
    return out
# 10. sorting & searching
sortOn = lambda f,xs: sorted(xs,key=f)
sortBy = lambda cmp,xs: sorted(xs,key=cmp_to_key(cmp))
maxOn = lambda f,xs: max(xs,key=f)
minOn = lambda f,xs: min(xs,key=f)
merge = lambda a,b: list(_hmerge(a,b))
# 11. strings
lines = str.splitlines
unlines = lambda ls: "".join(l + "\n" for l in ls)
words = str.split
unwords = " ".join
# 12. containers
def fromListWith(f,pairs):
    d = {}
    for k, v in pairs:
        d[k] = f(v,d[k]) if k in d else v
    return d
insertWith = lambda f,k,v,d: {**d, k: f(v,d[k]) if k in d else v}
unionWith = lambda f,a,b: {k: f(a[k],b[k]) if k in a and k in b else a[k] if k in a else b[k]
                           for k in a.keys() | b.keys()}
def getpath(t,ks,default=NOTHING):
    for k in ks:
        if not isinstance(t,dict) or k not in t: return default
        t = t[k]
    return t
bag_diff = lambda a,b: {k: a[k] - b.get(k,0) for k in a if a[k] > b.get(k,0)}
bag_inter = lambda a,b: {k: v for k in a.keys() & b.keys() if (v := min(a[k],b[k]))}
bag_union = lambda a,b: {k: max(a.get(k,0),b.get(k,0)) for k in a.keys() | b.keys()}
bag_sub = lambda a,b: all(b.get(k,0) >= n for k,n in a.items())
# 13. Maybe
bind = lambda x,f: NOTHING if x is NOTHING else f(x)
isJust = lambda x: x is not NOTHING
isNothing = lambda x: x is NOTHING
fromMaybe = lambda default,x: default if x is NOTHING else x
listToMaybe = lambda xs: xs[0] if xs else NOTHING
maybeToList = lambda x: [] if x is NOTHING else [x]
catMaybes = lambda xs: [x for x in xs if x is not NOTHING]
mapMaybe = lambda f,xs: [y for y in map(f,xs) if y is not NOTHING]
def sequenceM(xs):
    out = []
    for x in xs:
        if x is NOTHING: return NOTHING
        out.append(x)
    return out
traverseM = lambda f,xs: sequenceM(map_(f,xs))
maybe_get = lambda d,k: d.get(k,NOTHING)
# 14. Either
Ok = lambda v: ("ok", v)
Err = lambda e: ("err", e)
bindE = lambda x,f: x if x[0] == 'err' else f(x[1])
def sequenceE(xs):
    out = []
    for x in xs:
        if x[0] == "err": return x
        out.append(x[1])
    return Ok(out)
traverseE = lambda f,xs: sequenceE(map_(f,xs))
def partitionEithers(xs):
    oks, errs = [], []
    for tag, v in xs:
        (oks if tag == "ok" else errs).append(v)
    return (oks, errs)
note = lambda msg,x: Err(msg) if x is NOTHING else Ok(x)
# 15. do-notation
def do(binder,is_fail):
    def deco(gen_fn):
        def run(*a,**kw):
            g = gen_fn(*a,**kw)
            def step(val):
                try:
                    m = g.send(val)
                except StopIteration as e:
                    return e.value
                return binder(m,step)
            return step(None)
        return run
    return deco
doM = do(bind,lambda x: x is NOTHING)
doE = do(bindE,lambda x: x[0] == "err")
# 16. parser combinators
def doP(gen_fn):
    def make(*a,**kw):
        def parser(s):
            g, val = gen_fn(*a,**kw), None
            try:
                while True:
                    r = g.send(val)(s)
                    if r is FAIL: return FAIL
                    val, s = r
            except StopIteration as e:
                return (e.value, s)
        return parser
    return make
alt = lambda *ps: lambda s: next((r for r in (p(s) for p in ps) if r is not FAIL),FAIL)
def many(p):
    def parser(s):
        out = []
        while (r := p(s)) is not FAIL and r[1] != s:
            out.append(r[0]); s = r[1]
        return (out, s)
    return parser
def many1(p):
    def parser(s):
        r = p(s)
        if r is FAIL: return FAIL
        rest = many(p)(r[1])
        return ([r[0]] + rest[0], rest[1])
    return parser
def sepBy(p,sep):
    def pair(s):
        r = sep(s)
        return FAIL if r is FAIL else p(r[1])
    def parser(s):
        r = p(s)
        if r is FAIL: return ([], s)
        rest = many(pair)(r[1])
        return ([r[0]] + rest[0], rest[1])
    return parser
def lit(c):
    def parser(s):
        s = s.lstrip()
        return (c, s[1:]) if s[:1] == c else FAIL
    return parser
def rx(pat,conv=identity):
    r = _re.compile(pat)
    def parser(s):
        s = s.lstrip()
        m = r.match(s)
        return (conv(m.group()), s[m.end():]) if m else FAIL
    return parser
def chainl1(p,ops):
    opp = alt(*[lit(c) for c in ops])
    def parser(s):
        vr = p(s)
        while vr is not FAIL:
            op = opp(vr[1])
            if op is FAIL: return vr
            w = p(op[1])
            if w is FAIL: return vr
            vr = (ops[op[0]](vr[0],w[0]), w[1])
        return FAIL
    return parser
def runParser(p,s):
    r = p(s)
    if r is FAIL: return Err(f"no parse: {s!r}")
    if r[1].strip(): return Err(f"unconsumed input: {r[1].strip()!r}")
    return Ok(r[0])
# 17. monoids
Monoid = lambda empty,op: (empty, op)
mconcat = lambda m,xs: foldl(m[1],xs,m[0])
foldMap = lambda f,m,xs: mconcat(m,[f(x) for x in xs])
both = lambda m1,m2: Monoid((m1[0], m2[0]),
              lambda a,b: (m1[1](a[0],b[0]), m2[1](a[1],b[1])))
Sum = Monoid(0,add)
Product = Monoid(1,mul)
All = Monoid(True,lambda a,b: a and b)
Any = Monoid(False,lambda a,b: a or b)
MaxM = Monoid(float("-inf"),max)
MinM = Monoid(float("inf"),min)
ListM = Monoid([],add)
First = Monoid(NOTHING,lambda a,b: b if a is NOTHING else a)
Last = Monoid(NOTHING,lambda a,b: a if b is NOTHING else b)

\end{lstlisting}
\chapter{Point-Free Foundations}

A function that names its argument commits to one shape of input; a function
built from other functions (\texttt{compose}, \texttt{on}, \texttt{partial})
stays a transformation, reusable wherever the types line up. This chapter's
five one-liners (\texttt{identity}, \texttt{const}, \texttt{flip},
\texttt{curry}, \texttt{uncurry}) are the adapters that make composition
possible at all -- they change nothing about the values, only the SHAPE of
the call. Two sentinels close the chapter: \texttt{NOTHING} means ``no
answer'', and \texttt{FAIL} means ``the parser gave up'' -- both real
objects, never \texttt{None}, so a function can return an ordinary
\texttt{None} as genuine data without it being mistaken for absence.
\section{Homework}
Work in order: drills cement each tool, applies combine them,
challenges are interview-grade. Write each solution in a scratch file with
the given doctests pasted in, and let \texttt{python3 -m doctest} referee.
\subsection*{2.1\quad Mirror a pair\hfill\textnormal{\textit{drill}}}
Write mirror(p): the pair with its elements exchanged, point-free. Applying it twice must give the original back.

\noindent\texttt{mirror(p) -> pair}
\begin{lstlisting}[style=ex]
>>> mirror((1, 2))
(2, 1)
>>> mirror(mirror((3, 4)))
(3, 4)
\end{lstlisting}
\noindent\textbf{Solution.}
\begin{lstlisting}[style=code]
mirror = o(tuple, reversed)
\end{lstlisting}
\noindent\textit{Note.} mirror = o(tuple, reversed) is composition, read right to left: reversed(p) yields the pair's two elements back to front (like an iterator over (p[1], p[0])), and tuple(...) then materializes that into an actual tuple. Because the whole thing is built from function composition, no argument is ever named or indexed -- there is no p[0]/p[1] anywhere in the definition.

\subsection*{2.2\quad Record projections\hfill\textnormal{\textit{drill}}}
A leaderboard record is a (name, score) pair. Write name\_of and score\_of point-free, without lambdas.

\noindent\texttt{name\_of(r) -> value; score\_of(r) -> value}
\begin{lstlisting}[style=ex]
>>> name_of(('ada', 95))
'ada'
>>> score_of(('ada', 95))
95
\end{lstlisting}
\noindent\textbf{Solution.}
\begin{lstlisting}[style=code]
name_of, score_of = itemgetter(0), itemgetter(1)
\end{lstlisting}
\noindent\textit{Note.} operator.itemgetter(0) and itemgetter(1) build extractor functions without writing lambda r: r[0] by hand -- itemgetter is exactly Haskell's fst/snd generalized to any index. Because itemgetter returns a genuine function object rather than a lambda, it slots directly into map(name\_of, records) or sorted(records, key=score\_of) with no wrapping needed.

\subsection*{2.3\quad Default when missing\hfill\textnormal{\textit{drill}}}
Write port\_of(cfg): the 'port' entry of a config dict, defaulting to 8080 only when the key is absent. A key holding None is real data and must survive.

\noindent\texttt{port\_of(cfg) -> value}
\begin{lstlisting}[style=ex]
>>> port_of({'port': 9000})
9000
>>> port_of({})
8080
>>> port_of({'port': None}) is None
True
\end{lstlisting}
\noindent\textbf{Solution.}
\begin{lstlisting}[style=code]
port_of = lambda cfg: fromMaybe(8080, maybe_get(cfg, 'port'))
\end{lstlisting}
\noindent\textit{Note.} maybe\_get(cfg, 'port') looks up the key and returns the NOTHING sentinel specifically when the key is absent, never confusing that with a key that is present but stored as None. fromMaybe(8080, ...) only substitutes its default in the NOTHING case, so an explicitly stored None passes through untouched -- which is exactly why the third test (port set to None) must come back as None, not 8080.

\subsection*{2.4\quad Clamp to a range\hfill\textnormal{\textit{applied}}}
Write clamp(lo, hi): a FUNCTION that forces a number into [lo, hi]. Build it from max and min with partial application -- no lambda in the returned function.

\noindent\texttt{clamp(lo, hi) -> (x -> value)}
\begin{lstlisting}[style=ex]
>>> clamp(0, 10)(15)
10
>>> clamp(0, 10)(-5)
0
>>> clamp(0, 10)(7)
7
\end{lstlisting}
\noindent\textbf{Solution.}
\begin{lstlisting}[style=code]
clamp = lambda lo, hi: o(partial(max, lo), partial(min, hi))
\end{lstlisting}
\noindent\textit{Note.} partial(max, lo) and partial(min, hi) each freeze one argument, turning max/min into one-argument functions (Haskell-style sections). Composing them right to left with o means min(hi, x) runs FIRST, capping the value from above, and its result feeds into max(lo, ...), which floors it from below -- both bounds enforced in one pipeline, no if/elif needed.

\subsection*{2.5\quad Twice\hfill\textnormal{\textit{drill}}}
Write twice(f): the function that applies f two times.

\noindent\texttt{twice(f) -> (x -> value)}
\begin{lstlisting}[style=ex]
>>> twice(lambda n: n * 3)(2)
18
>>> twice(str.strip)('  hi  ')
'hi'
\end{lstlisting}
\noindent\textbf{Solution.}
\begin{lstlisting}[style=code]
twice = lambda f: o(f, f)
\end{lstlisting}
\noindent\textit{Note.} o(f, f) composes f with itself, so twice(f)(x) equals f(f(x)) by definition of composition. Because o is an ordinary function rather than special syntax, nothing stops it from taking the SAME function twice, and applying twice a second time -- twice(twice(f)) -- runs f four times, since composition itself composes.

\clearpage
\chapter{Folds, Scans and Unfolds}

A fold collapses a list to one value; a scan is a fold that keeps its
history; an unfold grows a list from a seed instead of consuming one.
\texttt{scanl} here is \emph{exactly} \texttt{itertools.accumulate} --
naming it the Haskell way does not change what runs. \texttt{unfoldr} is
the fold's mirror image: given a step function that returns
\texttt{(value, next\_seed)} or \texttt{NOTHING} to stop, it produces the
whole list, which is why balances, digit sequences and Collatz chains all
turn out to be one \texttt{unfoldr} call apiece. \texttt{iterate} is
\texttt{unfoldr}'s lazy, infinite cousin: it never stops on its own, so
\texttt{take} is what makes it usable.
\section{Homework}
Work in order: drills cement each tool, applies combine them,
challenges are interview-grade. Write each solution in a scratch file with
the given doctests pasted in, and let \texttt{python3 -m doctest} referee.
\subsection*{3.1\quad Parity census\hfill\textnormal{\textit{drill}}}
Count evens and odds as a pair (evens, odds), in ONE pass, point-free at the top level.

\noindent\texttt{parity\_counts(xs) -> (evens, odds)}
\begin{lstlisting}[style=ex]
>>> parity_counts([1, 2, 3, 4])
(2, 2)
>>> parity_counts([])
(0, 0)
>>> parity_counts([-2, -3])
(1, 1)
\end{lstlisting}
\noindent\textbf{Solution.}
\begin{lstlisting}[style=code]
parity_counts = partial(foldMap, lambda x: (1 - x % 2, x % 2), both(Sum, Sum))
\end{lstlisting}
\noindent\textit{Note.} foldMap first measures each element with lambda x: (1 - x \% 2, x \% 2), turning every value into a (is\_even, is\_odd) pair of 0s and 1s; both(Sum, Sum) then runs two independent Sum folds in lockstep, one over each tuple position. partial freezes the measuring function and the monoid together into a single reusable name, and because Sum's identity is 0, the empty-list case ((0, 0)) falls out for free with no special-casing.

\subsection*{3.2\quad Running balance\hfill\textnormal{\textit{drill}}}
An account starts at 0; deposits are positive, withdrawals negative. Return every intermediate balance including the start.

\noindent\texttt{balances(txs) -> [balance]}
\begin{lstlisting}[style=ex]
>>> balances([100, -30, -30, 5])
[0, 100, 70, 40, 45]
>>> balances([])
[0]
\end{lstlisting}
\noindent\textbf{Solution.}
\begin{lstlisting}[style=code]
balances = lambda txs: scanl(add, txs, 0)
\end{lstlisting}
\noindent\textit{Note.} scanl(add, txs, 0) is a fold that keeps EVERY intermediate accumulator instead of just the final one: starting from the seed 0, it emits the running total after each transaction in txs in turn, so the result is one element longer than the input (it includes the starting balance before any transaction is applied). operator.add stands in for lambda a, b: a + b, following the file's point-free convention.

\subsection*{3.3\quad Countdown\hfill\textnormal{\textit{drill}}}
Write countdown(n): the list n, n-1, ..., 1 grown by unfoldr, point-free at the top level.

\noindent\texttt{countdown(n) -> [int]}
\begin{lstlisting}[style=ex]
>>> countdown(4)
[4, 3, 2, 1]
>>> countdown(0)
[]
\end{lstlisting}
\noindent\textbf{Solution.}
\begin{lstlisting}[style=code]
countdown = partial(unfoldr, lambda n: NOTHING if n == 0 else (n, n - 1))
\end{lstlisting}
\noindent\textit{Note.} unfoldr grows a list from a seed by repeatedly calling its step function and collecting what it returns. Here the step returns (n, n - 1) -- this element, next seed -- for every n except 0, and returns NOTHING once n hits 0, which unfoldr treats as the stop signal. partial freezes the step function in, so countdown ends up a plain function of n alone.

\subsection*{3.4\quad Digits of a number\hfill\textnormal{\textit{applied}}}
Write digits(n): the base-10 digits of a non-negative integer, most significant first. Zero has one digit.

\noindent\texttt{digits(n) -> [int]}
\begin{lstlisting}[style=ex]
>>> digits(407)
[4, 0, 7]
>>> digits(0)
[0]
\end{lstlisting}
\noindent\textbf{Solution.}
\begin{lstlisting}[style=code]
digits = lambda n: unfoldr(lambda k: NOTHING if k == 0 else (k % 10, k // 10), n)[::-1] or [0]
\end{lstlisting}
\noindent\textit{Note.} The step function peels off the LEAST-significant digit first (k \% 10) and passes what remains onward as the next seed (k // 10), stopping at NOTHING once k reaches 0. Because digits emerge least-significant-first, [::-1] is needed to restore normal reading order; 'or [0]' exists because unfoldr(..., 0) immediately returns an empty list (0 stops right away), yet digits(0) must be [0], not [].

\subsection*{3.5\quad Collatz flight length\hfill\textnormal{\textit{challenge}}}
Write collatz\_len(n): how many numbers the Collatz sequence visits from n down to 1, inclusive. Even halves; odd triples-plus-one.

\noindent\texttt{collatz\_len(n) -> int}
\begin{lstlisting}[style=ex]
>>> collatz_len(6)
9
>>> collatz_len(1)
1
\end{lstlisting}
\noindent\textbf{Solution.}
\begin{lstlisting}[style=code]
collatz_len = lambda n: 1 + len(unfoldr(lambda k: NOTHING if k == 1 else (k, k // 2 if k % 2 == 0 else 3 * k + 1), n))
\end{lstlisting}
\noindent\textit{Note.} The step function encodes the whole rule in one line: halve k when it's even, otherwise 3k + 1, and stop (NOTHING) once k reaches 1. unfoldr therefore emits every number in the sequence EXCEPT the terminal 1, so len(...) counts those steps and +1 accounts for the 1 that was never emitted -- no mutable counter, no while loop, just an unfold, a length, and an addition.

\clearpage
\chapter{The List Gaps}

The largest chapter because it is the least glamorous: the ordinary list
and dict work that every program needs, named consistently instead of
reached for ad hoc. \texttt{windows} and \texttt{groupBy} both slide a view
across a sequence, but for different reasons -- one for a fixed-width
lookback (moving averages), one for runs of equal elements (RLE,
compression). \texttt{transpose} is ragged-safe on purpose: short rows
simply drop out of later columns instead of raising, so uneven data never
needs a special case. \texttt{fromListWith} and \texttt{unionWith} are the
two dict-building folds worth memorising -- almost every ``group these'' or
``merge these two counters'' problem is one of them with the right
combining function.
\section{Homework}
Work in order: drills cement each tool, applies combine them,
challenges are interview-grade. Write each solution in a scratch file with
the given doctests pasted in, and let \texttt{python3 -m doctest} referee.
\subsection*{4.1\quad Run-length encoding\hfill\textnormal{\textit{drill}}}
Compress a sequence into (element, count) pairs per maximal run.

\noindent\texttt{rle(xs) -> [(x, n)]}
\begin{lstlisting}[style=ex]
>>> rle('aaabbc')
[('a', 3), ('b', 2), ('c', 1)]
>>> rle([])
[]
\end{lstlisting}
\noindent\textbf{Solution.}
\begin{lstlisting}[style=code]
rle = lambda xs: [(g[0], len(g)) for g in groupBy(eq, xs)]
\end{lstlisting}
\noindent\textit{Note.} groupBy(eq, xs) partitions xs into consecutive runs of equal elements (operator.eq supplies the comparison, so no lambda is needed), returning something like [['a','a'],['b']]. The comprehension then turns each run g into (g[0], len(g)) -- the repeated value and how many times it repeated -- which is run-length encoding; groupBy has already done all the actual segmenting.

\subsection*{4.2\quad Moving average\hfill\textnormal{\textit{applied}}}
Write mavg(n, xs): the n-point moving average of a signal.

\noindent\texttt{mavg(n, xs) -> [float]}
\begin{lstlisting}[style=ex]
>>> mavg(3, [1, 2, 3, 4, 5])
[2.0, 3.0, 4.0]
>>> mavg(3, [1, 2])
[]
\end{lstlisting}
\noindent\textbf{Solution.}
\begin{lstlisting}[style=code]
mavg = lambda n, xs: [sum(w) / n for w in windows(n, xs)]
\end{lstlisting}
\noindent\textit{Note.} windows(n, xs) produces every contiguous length-n slice as its own list -- a genuine sliding window, not index arithmetic -- and the comprehension averages each one with sum(w) / n. Because windows naturally yields zero windows when xs is shorter than n, mavg returns [] for too-short input automatically, with no explicit length check anywhere.

\subsection*{4.3\quad Pagination\hfill\textnormal{\textit{drill}}}
Write paginate(n, xs): pages of at most n items, last page short.

\noindent\texttt{paginate(n, xs) -> [[x]]}
\begin{lstlisting}[style=ex]
>>> paginate(2, [1, 2, 3, 4, 5])
[[1, 2], [3, 4], [5]]
\end{lstlisting}
\noindent\textbf{Solution.}
\begin{lstlisting}[style=code]
paginate = chunksOf
\end{lstlisting}
\noindent\textit{Note.} This is a pure recognition exercise: chunksOf(n, xs) already means ''split xs into pages of size n'', with the final page possibly shorter -- there is nothing left to implement. paginate is chunksOf under a different name for a different domain.

\subsection*{4.4\quad Column sums, ragged\hfill\textnormal{\textit{applied}}}
Sensor rows arrive with trailing dropouts. Sum each column over the rows that reached it -- point-free.

\noindent\texttt{colsums(rows) -> [sum]}
\begin{lstlisting}[style=ex]
>>> colsums([[1, 2, 3], [4, 5]])
[5, 7, 3]
>>> colsums([])
[]
\end{lstlisting}
\noindent\textbf{Solution.}
\begin{lstlisting}[style=code]
colsums = o(list, partial(map, sum), transpose)
\end{lstlisting}
\noindent\textit{Note.} transpose(rows) turns rows into columns, and because it's ragged-safe, a row shorter than the others simply drops out of the columns past its own length instead of raising an error. Once the input is transposed, map(sum, ...) sums each resulting column, and o composes transpose, the mapping, and list-materialization right to left into one pipeline.

\subsection*{4.5\quad Indentation split\hfill\textnormal{\textit{drill}}}
Split a line into (leading whitespace, rest) as strings, one pass.

\noindent\texttt{indent\_split(s) -> (str, str)}
\begin{lstlisting}[style=ex]
>>> indent_split('   x = 1')
('   ', 'x = 1')
>>> indent_split('top')
('', 'top')
\end{lstlisting}
\noindent\textbf{Solution.}
\begin{lstlisting}[style=code]
indent_split = lambda s: tuple(''.join(p) for p in span(str.isspace, s))
\end{lstlisting}
\noindent\textit{Note.} span(str.isspace, s) walks s and splits it into (leading run of whitespace characters, everything after) as two lists of individual characters, stopping at the first non-whitespace character it meets. Because span works over s as a sequence of characters rather than as a string, ''.join(...) is needed on each half to turn the character lists back into real strings.

\clearpage
\chapter{Maybe, Done Right}

\texttt{None} is Python's ``no value'', which is exactly the problem: a
function cannot then return \texttt{None} to mean a REAL answer without
creating an ambiguity. \texttt{NOTHING} fixes this by being a distinct
object that is never a legitimate value -- \texttt{maybe\_get} returns
\texttt{None} for a key whose stored value truly is \texttt{None}, and
\texttt{NOTHING} only when the key is missing outright. \texttt{bind}
chains failable steps by short-circuiting the moment one returns
\texttt{NOTHING}; \texttt{sequenceM} and \texttt{traverseM} extend that to
a whole list, all-or-nothing. Once a pipeline's failure mode is a value
instead of an exception, the pipeline composes like any other function.
\section{Homework}
Work in order: drills cement each tool, applies combine them,
challenges are interview-grade. Write each solution in a scratch file with
the given doctests pasted in, and let \texttt{python3 -m doctest} referee.
\subsection*{5.1\quad Safe head\hfill\textnormal{\textit{drill}}}
Write headM(xs): the first element, or NOTHING for an empty sequence. A leading None is data.

\noindent\texttt{headM(xs) -> Maybe value}
\begin{lstlisting}[style=ex]
>>> headM([7, 8])
7
>>> headM([]) is NOTHING
True
>>> headM([None]) is None
True
\end{lstlisting}
\noindent\textbf{Solution.}
\begin{lstlisting}[style=code]
headM = lambda xs: xs[0] if xs else NOTHING
\end{lstlisting}
\noindent\textit{Note.} The whole exercise is what survives the check: 'xs[0] if xs else NOTHING' returns the first element whenever the list is non-empty -- even if that element genuinely IS None, a real stored value -- and only returns the NOTHING sentinel when the list is truly empty. Using plain None to mean ''empty'' would make a list starting with None indistinguishable from an empty list, which is exactly the ambiguity the third test is built to catch.

\subsection*{5.2\quad Deep dictionary path\hfill\textnormal{\textit{applied}}}
Write pathM(d, ks): follow a key path through nested dicts; NOTHING if any hop is missing. A stored None must survive the trip.

\noindent\texttt{pathM(d, ks) -> Maybe value}
\begin{lstlisting}[style=ex]
>>> d = {'db': {'host': 'local', 'pool': None}}
>>> pathM(d, ['db', 'host'])
'local'
>>> pathM(d, ['db', 'pool']) is None
True
>>> pathM(d, ['db', 'user']) is NOTHING
True
\end{lstlisting}
\noindent\textbf{Solution.}
\begin{lstlisting}[style=code]
pathM = lambda d, ks: foldl(lambda m, k: bind(m, lambda t: maybe_get(t, k)), ks, d)
\end{lstlisting}
\noindent\textit{Note.} foldl threads the dictionary through each key in ks one at a time; at every step, bind(m, lambda t: maybe\_get(t, k)) either descends into the next nested dict (if the lookup so far succeeded) or short-circuits straight to NOTHING the moment any key along the path is missing. This is the Maybe monad's signature move: a chain of possibly-failing steps where one failure propagates to the very end without an if-chain anywhere.

\subsection*{5.3\quad Salvage the parseable\hfill\textnormal{\textit{drill}}}
Given strings, keep the ones that are integers, as ints, dropping the rest.

\noindent\texttt{salvage\_ints(ss) -> [int]}
\begin{lstlisting}[style=ex]
>>> salvage_ints(['3', 'x', '-7', ''])
[3, -7]
\end{lstlisting}
\noindent\textbf{Solution.}
\begin{lstlisting}[style=code]
salvage_ints = lambda ss: catMaybes([int(s) if s.lstrip('-').isdigit() and s else NOTHING for s in ss])
\end{lstlisting}
\noindent\textit{Note.} Each string passes through an inline test-then-convert expression that yields either its integer value or NOTHING (rejecting empty strings and anything that isn't a valid, possibly-negative integer); catMaybes then filters out every NOTHING, keeping only the strings that actually parsed. This is mapMaybe's job split into two explicit, named steps.

\subsection*{5.4\quad All keys or nothing\hfill\textnormal{\textit{applied}}}
Write need(d, ks): the values for ks in order, or NOTHING if ANY key is absent.

\noindent\texttt{need(d, ks) -> Maybe [value]}
\begin{lstlisting}[style=ex]
>>> need({'a': 1, 'b': 2}, ['a', 'b'])
[1, 2]
>>> need({'a': 1}, ['a', 'b']) is NOTHING
True
\end{lstlisting}
\noindent\textbf{Solution.}
\begin{lstlisting}[style=code]
need = lambda d, ks: sequenceM([maybe_get(d, k) for k in ks])
\end{lstlisting}
\noindent\textit{Note.} sequenceM is handed a list of Maybe values, one maybe\_get per requested key, and returns the fully unwrapped list only if EVERY lookup succeeded; the instant any single lookup comes back NOTHING, the whole result collapses to NOTHING. That is precisely ''all keys present, or nothing at all'' -- no manual all()-style checking required.

\subsection*{5.5\quad First responder\hfill\textnormal{\textit{challenge}}}
Several config sources; write first\_set(dicts, k): the value from the FIRST dict that defines k at all -- even as None -- else NOTHING.

\noindent\texttt{first\_set(dicts, k) -> Maybe value}
\begin{lstlisting}[style=ex]
>>> first_set([{}, {'x': None}, {'x': 3}], 'x') is None
True
>>> first_set([{}, {}], 'x') is NOTHING
True
\end{lstlisting}
\noindent\textbf{Solution.}
\begin{lstlisting}[style=code]
first_set = lambda dicts, k: mconcat(First, [maybe_get(d, k) for d in dicts])
\end{lstlisting}
\noindent\textit{Note.} mconcat folds the list of per-dict lookups through the First monoid, whose combining rule is ''keep the first non-NOTHING value seen, ignore the rest''; because First's identity element IS NOTHING, an empty list of dicts (or one where every dict is missing the key) correctly yields NOTHING with no base case to write by hand.

\clearpage
\chapter{Either: Failure With a Why}

Maybe answers ``did it work?''; Either answers that AND ``why not?''.
\texttt{Ok}/\texttt{Err} are tagged pairs -- \texttt{("ok", v)} or
\texttt{("err", why)} -- so \texttt{bindE} can short-circuit exactly like
\texttt{bind} while still carrying the failure's reason to wherever the
pipeline is finally inspected. \texttt{sequenceE} and \texttt{traverseE}
are Either's all-or-nothing gates, and \texttt{note} is the bridge from the
last chapter: it turns a \texttt{NOTHING} into an \texttt{Err} with a
message attached, for the moment a Maybe pipeline needs to explain itself
to a caller.
\section{Homework}
Work in order: drills cement each tool, applies combine them,
challenges are interview-grade. Write each solution in a scratch file with
the given doctests pasted in, and let \texttt{python3 -m doctest} referee.
\subsection*{6.1\quad Lift the miss into a reason\hfill\textnormal{\textit{drill}}}
Write get\_e(d, k): Ok(value) or Err naming the missing key.

\noindent\texttt{get\_e(d, k) -> Either}
\begin{lstlisting}[style=ex]
>>> get_e({'a': 1}, 'a')
('ok', 1)
>>> get_e({}, 'user')
('err', "missing key: 'user'")
\end{lstlisting}
\noindent\textbf{Solution.}
\begin{lstlisting}[style=code]
get_e = lambda d, k: note(f"missing key: {k!r}", maybe_get(d, k))
\end{lstlisting}
\noindent\textit{Note.} note is the bridge from Maybe to Either: given a message and a Maybe value, it turns a NOTHING into Err(message) and leaves any real value -- including a successful None -- wrapped as Ok(value), unchanged. A lookup that could previously only say ''missing'' can now also say WHY.

\subsection*{6.2\quad Parse then bound\hfill\textnormal{\textit{applied}}}
Write to\_port(s): parse a decimal string and demand 1..65535, each failure carrying its own message.

\noindent\texttt{to\_port(s) -> Either}
\begin{lstlisting}[style=ex]
>>> to_port('8080')
('ok', 8080)
>>> to_port('http')
('err', "not a number: 'http'")
>>> to_port('99999')
('err', 'out of range: 99999')
\end{lstlisting}
\noindent\textbf{Solution.}
\begin{lstlisting}[style=code]
to_port = lambda s: bindE(
    note(f"not a number: {s!r}", int(s) if s.isdigit() else NOTHING),
    lambda n: Ok(n) if 1 <= n <= 65535 else Err(f"out of range: {n}"))
\end{lstlisting}
\noindent\textit{Note.} bindE chains two Either-producing steps. First, note converts a failed int() parse (guarded by isdigit()) into an Err carrying a parsing message; if that succeeds, the second step (the lambda) checks the number falls in the valid port range and returns either Ok(n) or a DIFFERENT Err for an out-of-range value. Both failure reasons flow through the same bindE plumbing, so the caller inspects one Either at the end instead of two separate failure sites.

\subsection*{6.3\quad First error speaks\hfill\textnormal{\textit{drill}}}
Validate a batch: all Ok values, or the FIRST Err verbatim.

\noindent\texttt{all\_ports(ss) -> Either}
\begin{lstlisting}[style=ex]
>>> all_ports(['80', '443'])
('ok', [80, 443])
>>> all_ports(['80', 'x', 'y'])
('err', "not a number: 'x'")
\end{lstlisting}
\noindent\textbf{Solution.}
\begin{lstlisting}[style=code]
all_ports = lambda ss: sequenceE([to_port(s) for s in ss])
\end{lstlisting}
\noindent\textit{Note.} sequenceE applies the same all-or-nothing discipline as sequenceM, but for Either: it walks the list of per-string results from to\_port and returns Ok of the whole list only if every one succeeded, otherwise it returns the FIRST Err it hit, message intact. Unlike sequenceM's bare NOTHING, the caller learns exactly which port was invalid and why.

\subsection*{6.4\quad Division with a why\hfill\textnormal{\textit{drill}}}
Write div\_e(a, b): the quotient, or an error naming the offense.

\noindent\texttt{div\_e(a, b) -> Either}
\begin{lstlisting}[style=ex]
>>> div_e(10, 4)
('ok', 2.5)
>>> div_e(1, 0)
('err', 'division by zero')
\end{lstlisting}
\noindent\textbf{Solution.}
\begin{lstlisting}[style=code]
div_e = lambda a, b: Err('division by zero') if b == 0 else Ok(a / b)
\end{lstlisting}
\noindent\textit{Note.} This is the base case the rest of the chapter's Either machinery is built on: a genuinely failable operation that returns an Err with an explanation instead of raising ZeroDivisionError, so the failure can be caught, inspected, and composed with bindE/doE exactly like every other Either-returning step in this chapter.

\clearpage
\chapter{do-Notation}

Chaining \texttt{bind} calls by hand nests one call inside the next, one
level per step -- the ``bind pyramid''. \texttt{do} rebuilds that same
chain out of an ordinary generator function: each \texttt{yield} is one
\texttt{bind}, and the value sent back in is what the step would have
returned. \texttt{doM} wires this to the Maybe monad, \texttt{doE} to
Either -- same decorator, same generator shape, different short-circuit
rule. The win is not cleverness for its own sake: a five-step Maybe
pipeline written with \texttt{doM} reads top to bottom, in the order the
steps actually happen, with the failure path made invisible instead of
handled at every line.
\section{Homework}
Work in order: drills cement each tool, applies combine them,
challenges are interview-grade. Write each solution in a scratch file with
the given doctests pasted in, and let \texttt{python3 -m doctest} referee.
\subsection*{7.1\quad Config reader, flat\hfill\textnormal{\textit{drill}}}
Using doM, read host and port from a dict; NOTHING if either is absent.

\noindent\texttt{hostport(d) -> Maybe (host, port)}
\begin{lstlisting}[style=ex]
>>> hostport({'host': 'h', 'port': 80})
('h', 80)
>>> hostport({'host': 'h'}) is NOTHING
True
\end{lstlisting}
\noindent\textbf{Solution.}
\begin{lstlisting}[style=code]
@doM
def hostport(d):
    h = yield maybe_get(d, 'host')
    p = yield maybe_get(d, 'port')
    return (h, p)
\end{lstlisting}
\noindent\textit{Note.} @doM turns an ordinary generator into a Maybe pipeline: each yield hands a Maybe value to doM's machinery, which unwraps it and sends the unwrapped value back into the generator as the result of that yield expression if it's real, or aborts the WHOLE function with NOTHING the instant any yield produces NOTHING. Written this way, hostport reads as two straight-line lookups with no visible branching, even though either one could fail.

\subsection*{7.2\quad Join across two tables\hfill\textnormal{\textit{applied}}}
users maps id to name; emails maps name to address. Using doM, resolve an id to its address.

\noindent\texttt{address(users, emails, uid) -> Maybe str}
\begin{lstlisting}[style=ex]
>>> u, e = {1: 'ada'}, {'ada': 'ada@x.io'}
>>> address(u, e, 1)
'ada@x.io'
>>> address(u, e, 2) is NOTHING
True
\end{lstlisting}
\noindent\textbf{Solution.}
\begin{lstlisting}[style=code]
@doM
def address(users, emails, uid):
    name = yield maybe_get(users, uid)
    return (yield maybe_get(emails, name))
\end{lstlisting}
\noindent\textit{Note.} The same @doM discipline, but the second lookup's key depends on the FIRST lookup's result: name = yield maybe\_get(users, uid) must succeed and bind name before maybe\_get(emails, name) can even be attempted. If uid isn't in users, the function stops at the first yield and never reaches the second lookup -- there is nothing to explicitly check for that.

\subsection*{7.3\quad Validated record\hfill\textnormal{\textit{applied}}}
Using doE, build a user record from a form dict: name must be non-empty, age must parse and land in 0..130. Each rejection explains itself.

\noindent\texttt{mk\_user(form) -> Either dict}
\begin{lstlisting}[style=ex]
>>> mk_user({'name': 'ada', 'age': '36'})
('ok', {'name': 'ada', 'age': 36})
>>> mk_user({'name': '', 'age': '36'})
('err', 'empty name')
>>> mk_user({'name': 'ada', 'age': 'old'})
('err', "bad age: 'old'")
\end{lstlisting}
\noindent\textbf{Solution.}
\begin{lstlisting}[style=code]
@doE
def mk_user(form):
    name = yield (Ok(form.get('name', '')) if form.get('name') else Err('empty name'))
    raw  = form.get('age', '')
    age  = yield note(f"bad age: {raw!r}", int(raw) if raw.isdigit() else NOTHING)
    yield (Ok(age) if age <= 130 else Err(f'implausible age: {age}'))
    return Ok({'name': name, 'age': age})
\end{lstlisting}
\noindent\textit{Note.} This doE pipeline mixes VALUE-producing yields (which bind name and age for later use) with one bare, unbound yield -- the age <= 130 check -- that runs purely for its side effect of stopping the pipeline if it fails; nothing from that check is kept, it exists only to guard what comes after. Read top to bottom, the happy path (every field valid) is the only path visible, and every rejection reason sits as an Err right at the line that would fail it.

\subsection*{7.4\quad Chained arithmetic\hfill\textnormal{\textit{drill}}}
Using doE and div\_e from 6.4, compute a/b + c/d with the first bad divisor reported.

\noindent\texttt{sum\_of\_ratios(a, b, c, d) -> Either}
\begin{lstlisting}[style=ex]
>>> sum_of_ratios(1, 2, 1, 4)
('ok', 0.75)
>>> sum_of_ratios(1, 0, 1, 4)
('err', 'division by zero')
\end{lstlisting}
\noindent\textbf{Solution.}
\begin{lstlisting}[style=code]
@doE
def sum_of_ratios(a, b, c, d):
    x = yield div_e(a, b)
    y = yield div_e(c, d)
    return Ok(x + y)
\end{lstlisting}
\noindent\textit{Note.} Two divisions, each already Either-safe via div\_e, are bound in sequence with doE: x captures the first ratio, y the second, and if EITHER division is by zero, the whole function short-circuits to that division's Err without ever reaching the addition. The final line is pure arithmetic with no visible error-handling code, because doE has already absorbed all of it.

\clearpage
\chapter{Parser Combinators}

A parser is a function from a string to \texttt{(value, rest)} or
\texttt{FAIL} -- once that shape is fixed, parsers compose like any other
function. \texttt{doP} threads the remaining input through a generator the
same way \texttt{doM} threads a Maybe; \texttt{alt} tries alternatives in
order; \texttt{many} and \texttt{sepBy} handle repetition and
separator-delimited lists without recursion. \texttt{chainl1} deserves
special attention: it folds a sequence of \texttt{p (op p)*} LEFT, which is
exactly how left-associative arithmetic is supposed to parse, and it is the
one combinator in this chapter doing real algorithmic work rather than
just gluing others together.
\section{Homework}
Work in order: drills cement each tool, applies combine them,
challenges are interview-grade. Write each solution in a scratch file with
the given doctests pasted in, and let \texttt{python3 -m doctest} referee.
\subsection*{8.1\quad Comma-separated integers\hfill\textnormal{\textit{drill}}}
Parse a comma-separated integer list, whitespace-blind, whole input consumed.

\noindent\texttt{int\_list(s) -> Either [int]}
\begin{lstlisting}[style=ex]
>>> int_list('1, -2 , 3')
('ok', [1, -2, 3])
>>> int_list('1, x')
('err', "unconsumed input: ', x'")
\end{lstlisting}
\noindent\textbf{Solution.}
\begin{lstlisting}[style=code]
int_list = lambda s: runParser(sepBy(rx(r'-?\d+', int), lit(',')), s)
\end{lstlisting}
\noindent\textit{Note.} sepBy(rx(...), lit(',')) parses one or more integer tokens separated by literal commas, correctly handling the fencepost problem (n items need only n - 1 separators) with no manual counting. runParser then additionally demands the ENTIRE input be consumed (rejecting trailing garbage) and reports the outcome as an Either instead of a bare value-or-FAIL.

\subsection*{8.2\quad key=value line\hfill\textnormal{\textit{drill}}}
Parse one 'key = value' assignment into a pair; keys are words, values run to end of line.

\noindent\texttt{kv(s) -> Either (k, v)}
\begin{lstlisting}[style=ex]
>>> kv('host = 10.0.0.1')
('ok', ('host', '10.0.0.1'))
>>> kv('= oops')[0]
'err'
\end{lstlisting}
\noindent\textbf{Solution.}
\begin{lstlisting}[style=code]
@doP
def _kv():
    k = yield rx(r'\w+')
    yield lit('=')
    v = yield rx(r'[^\n]+', str.strip)
    return (k, v)
kvline = _kv()
kv = lambda s: runParser(kvline, s)
\end{lstlisting}
\noindent\textit{Note.} Three sequential yields inside a doP block: match an identifier, match and discard a literal '=', then match everything up to end of line (stripped) as the value. @doP threads the unconsumed remainder of the string from one yield to the next automatically, so no explicit ''rest of string'' variable is passed by hand anywhere. Defining kvline once means the same parser value can be reused later -- chapter 9's INI reader is built directly on it.

\subsection*{8.3\quad CSV row\hfill\textnormal{\textit{drill}}}
Split one unquoted CSV row into stripped fields.

\noindent\texttt{csv\_row(s) -> Either [str]}
\begin{lstlisting}[style=ex]
>>> csv_row('ada, 95, pass')
('ok', ['ada', '95', 'pass'])
\end{lstlisting}
\noindent\textbf{Solution.}
\begin{lstlisting}[style=code]
csv_row = lambda s: runParser(sepBy(rx(r'[^,\n]+', str.strip), lit(',')), s)
\end{lstlisting}
\noindent\textit{Note.} Structurally identical to 8.1's comma-separated integers, just with a token parser that matches and strips arbitrary non-comma text instead of digits. Recognizing that a CSV row and a list of integers are the same shape with a different token is the entire exercise -- sepBy itself doesn't change at all.

\subsection*{8.4\quad ISO date\hfill\textnormal{\textit{applied}}}
Parse YYYY-MM-DD into an (y, m, d) integer triple, rejecting month 13.

\noindent\texttt{iso\_date(s) -> Either (y, m, d)}
\begin{lstlisting}[style=ex]
>>> iso_date('2026-09-17')
('ok', (2026, 9, 17))
>>> iso_date('2026-13-01')
('err', 'bad month: 13')
\end{lstlisting}
\noindent\textbf{Solution.}
\begin{lstlisting}[style=code]
num = lambda n: rx(r'\d{%d}' % n, int)
@doP
def _date():
    y = yield num(4)
    yield lit('-')
    m = yield num(2)
    yield lit('-')
    d = yield num(2)
    return (y, m, d)
iso_date = lambda s: bindE(runParser(_date(), s),
    lambda t: Ok(t) if 1 <= t[1] <= 12 else Err(f'bad month: {t[1]}'))
\end{lstlisting}
\noindent\textit{Note.} \_date parses three fixed-width numeric fields (4 digits, 2 digits, 2 digits) separated by literal dashes via doP, producing a raw (year, month, day) tuple purely from the STRING SHAPE, with no awareness of what a valid month even is. bindE then takes that successfully-parsed tuple and applies a semantic check -- is the month between 1 and 12 -- as a separate stage, turning an otherwise-valid parse into an Err if the calendar rule is broken. Keeping ''does this look like a date'' and ''is this a real date'' as two stages means either can change without touching the other.

\subsection*{8.5\quad Sums and differences\hfill\textnormal{\textit{challenge}}}
Parse and evaluate a left-associative +/- chain over integers.

\noindent\texttt{sums(s) -> Either int}
\begin{lstlisting}[style=ex]
>>> sums('10 - 3 - 2 + 1')
('ok', 6)
>>> sums('10 -')[0]
'err'
\end{lstlisting}
\noindent\textbf{Solution.}
\begin{lstlisting}[style=code]
sums = lambda s: runParser(chainl1(rx(r'-?\d+', int), {'+': add, '-': sub}), s)
\end{lstlisting}
\noindent\textit{Note.} chainl1 takes a token parser (digits) and a dict mapping operator characters to their Python functions, and folds a sequence of ''number, operator, number, operator, number...'' strictly left to right, matching how ordinary arithmetic associates addition and subtraction. operator.add/sub are used directly as the combining functions instead of hand-written lambdas, keeping the whole solution point-free.

\clearpage
\chapter{Monoids at Work}

A monoid is nothing but an identity element paired with an associative
combiner, reified as the pair \texttt{(empty, op)} -- naming it as DATA
means one engine, \texttt{mconcat}, folds every instance, and
\texttt{foldMap} fuses ``measure each element'' with ``combine the
measurements'' into one pass. \texttt{both} is the chapter's real payoff:
it pairs two monoids into one, so a max and a min, or a count and a total,
fall out of a SINGLE traversal instead of two separate loops. That matters
most exactly when it looks like it should not -- a one-shot iterator or a
huge file that can only be read once.
\section{Homework}
Work in order: drills cement each tool, applies combine them,
challenges are interview-grade. Write each solution in a scratch file with
the given doctests pasted in, and let \texttt{python3 -m doctest} referee.
\subsection*{9.1\quad Min and max, one pass\hfill\textnormal{\textit{drill}}}
Both extremes of a numeric stream in a single traversal, point-free.

\noindent\texttt{extremes(xs) -> (min, max)}
\begin{lstlisting}[style=ex]
>>> extremes([3, 1, 4, 1, 5])
(1, 5)
\end{lstlisting}
\noindent\textbf{Solution.}
\begin{lstlisting}[style=code]
extremes = partial(foldMap, lambda x: (x, x), both(MinM, MaxM))
\end{lstlisting}
\noindent\textit{Note.} foldMap first turns each element x into the pair (x, x) -- one copy destined for the min side, one for the max side -- and both(MinM, MaxM) runs MinM's fold on the first half of every pair and MaxM's fold on the second half, all in a SINGLE traversal of xs. That matters specifically when xs is a generator or a large file read once: two separate min(xs)/max(xs) calls would each need their own full pass, impossible for a one-shot iterator, but this needs only one.

\subsection*{9.2\quad Stream statistics\hfill\textnormal{\textit{applied}}}
Count, total and maximum of a stream in ONE pass: ((count, total), maximum).

\noindent\texttt{stats(xs) -> ((n, total), mx)}
\begin{lstlisting}[style=ex]
>>> stats([3, 1, 4])
((3, 8), 4)
\end{lstlisting}
\noindent\textbf{Solution.}
\begin{lstlisting}[style=code]
stats = partial(foldMap, lambda x: ((1, x), x), both(both(Sum, Sum), MaxM))
\end{lstlisting}
\noindent\textit{Note.} The same one-pass trick as 8.1, nested one level deeper: each element becomes ((1, x), x), and both(both(Sum, Sum), MaxM) applies Sum to the running count, Sum to the running total, and MaxM to the running maximum, all three statistics falling out of three nested monoids fused into a single fold. The output shape, ((n, total), mx), directly mirrors how the monoids were nested inside each other.

\subsection*{9.3\quad Config layering\hfill\textnormal{\textit{drill}}}
defaults, file, cli -- highest layer that DEFINES a key wins, and defining it as None counts. Point-free over the layer list.

\noindent\texttt{effective(layers, k) -> Maybe value}
\begin{lstlisting}[style=ex]
>>> layers = [{'v': 1}, {}, {'v': None}]
>>> effective(layers, 'v') is None
True
>>> effective(layers, 'x') is NOTHING
True
\end{lstlisting}
\noindent\textbf{Solution.}
\begin{lstlisting}[style=code]
effective = lambda layers, k: mconcat(Last, [maybe_get(d, k) for d in layers])
\end{lstlisting}
\noindent\textit{Note.} mconcat folds the list of per-layer lookups through the Last monoid, whose rule -- ''the most recent non-NOTHING value wins'' -- is precisely ''later layers override earlier ones''. Because Last's identity is NOTHING, a key present in no layer at all correctly yields NOTHING with no base case to write. This is the exact same skeleton as first\_set in 5.5, with only the monoid swapped (Last instead of First) to flip which end of the list wins.

\subsection*{9.4\quad Flatten\hfill\textnormal{\textit{drill}}}
Concatenate a list of lists by monoid fold.

\noindent\texttt{flatten(xss) -> [x]}
\begin{lstlisting}[style=ex]
>>> flatten([[1], [], [2, 3]])
[1, 2, 3]
\end{lstlisting}
\noindent\textbf{Solution.}
\begin{lstlisting}[style=code]
flatten = partial(mconcat, ListM)
\end{lstlisting}
\noindent\textit{Note.} mconcat(ListM, xss) folds the list of lists through ListM, the monoid whose identity is [] and whose op is list concatenation -- which is precisely what ''flatten a list of lists'' means. partial freezes ListM in, so flatten ends up a plain one-argument function. The quadratic-cost warning matters because repeated concatenation re-copies earlier elements on every combine; fine for a handful of sublists, but itertools.chain is the right tool once there are many large ones.

\clearpage
\chapter{Capstone Labs}

No new vocabulary -- this chapter is the payoff for the previous eight.
Each problem reaches back for tools from more than one chapter at once:
parsing with \texttt{doP}, then judging the parsed value with
\texttt{bindE}; folding a monoid across a batch validated with
\texttt{doE}. If a chapter's tools felt like isolated tricks on first
reading, this is where they stop being isolated.
\section{Homework}
Work in order: drills cement each tool, applies combine them,
challenges are interview-grade. Write each solution in a scratch file with
the given doctests pasted in, and let \texttt{python3 -m doctest} referee.
\subsection*{10.1\quad The algebraic calculator\hfill\textnormal{\textit{lab}}}
Full expression evaluator: + - * /, precedence, LEFT associativity, unary minus, parentheses, decimals, whitespace; Either out.

\noindent\texttt{calc(s) -> Either float}
\begin{lstlisting}[style=ex]
>>> calc('2 * (3 + 4) - -5')
('ok', 19.0)
>>> calc('8 - 3 - 2')
('ok', 3.0)
>>> calc('2 * (3 +')[0]
'err'
\end{lstlisting}
\noindent\textbf{Solution.}
\begin{lstlisting}[style=code]
number = rx(r'-?\d+(\.\d+)?', float)
@doP
def _paren():
    yield lit('(')
    v = yield expr
    yield lit(')')
    return v
unary = lambda s: alt(number, _paren())(s)
term  = chainl1(unary, {'*': mul, '/': lambda a, b: a / b})
expr  = chainl1(term,  {'+': add, '-': sub})
calc  = partial(runParser, expr)
\end{lstlisting}
\noindent\textit{Note.} A small recursive-descent grammar in three layers: unary handles a bare number OR a fully parenthesized sub-expression (recursing back into expr for whatever is inside the parens), term folds a sequence of unary results left to right through * and /, and expr folds a sequence of term results left to right through + and - -- together enforcing standard precedence purely through which layer calls which. doP threads the remaining input through every yield with no explicit position variable, chainl1 does the left-associative folding, and runParser turns the result into an Either that names exactly where parsing failed. Because none of this recurses per TOKEN (only per parenthesis nesting level), a long flat expression with thousands of terms still runs in constant stack depth.

\subsection*{10.2\quad INI reader\hfill\textnormal{\textit{lab}}}
Parse [section] headers with key = value bodies into a dict of dicts.

\noindent\texttt{ini(s) -> Either {section: {k: v}}}
\begin{lstlisting}[style=ex]
>>> t = '[db]\nhost = local\nport = 5432\n[app]\ndebug = true'
>>> ini(t)
('ok', {'db': {'host': 'local', 'port': '5432'}, 'app': {'debug': 'true'}})
\end{lstlisting}
\noindent\textbf{Solution.}
\begin{lstlisting}[style=code]
@doP
def _kvi():
    k = yield rx(r'\w+')
    yield lit('=')
    v = yield rx(r'[^\n\[]+', str.strip)
    return (k, v)
@doP
def _section():
    yield lit('[')
    name = yield rx(r'\w+')
    yield lit(']')
    kvs = yield many(_kvi())
    return (name, dict(kvs))
ini = lambda s: bindE(runParser(many(_section()), s), o(Ok, dict))
\end{lstlisting}
\noindent\textit{Note.} \_kvi parses one key=value line, reusing the exact shape built in 8.2, and \_section parses a bracketed header followed by ZERO OR MORE key/value lines via many(\_kvi()); many loops rather than recurses, so an arbitrarily long -- or empty -- section body costs no extra stack. The top-level ini parser is many(\_section()) run through runParser, and o(Ok, dict) takes the raw list of (name, {k: v}) pairs on success and wraps it as {name: {k: v}} inside an Ok, all in one composed step.

\subsection*{10.3\quad Log triage\hfill\textnormal{\textit{lab}}}
Lines are 'LEVEL message'. In ONE pass report ((errors, total), first ERROR message or NOTHING).

\noindent\texttt{triage(lines) -> ((errors, total), first\_err)}
\begin{lstlisting}[style=ex]
>>> triage(['INFO up', 'ERROR db down', 'ERROR retry', 'INFO ok'])
((2, 4), 'db down')
>>> triage(['INFO up'])[1] is NOTHING
True
\end{lstlisting}
\noindent\textbf{Solution.}
\begin{lstlisting}[style=code]
tag = lambda ln: (lambda lvl, _, msg: ((lvl == 'ERROR', 1),
                  msg if lvl == 'ERROR' else NOTHING))(*ln.partition(' '))
triage = partial(foldMap, tag, both(both(Sum, Sum), First))
\end{lstlisting}
\noindent\textit{Note.} tag turns each log line into a nested tuple shaped exactly like the monoid that will consume it: ((is\_it\_an\_error, 1), the\_error\_message\_or\_NOTHING). both(both(Sum, Sum), First) then counts errors, counts total lines, and remembers the first error's message, all as three independent monoids riding the SAME single pass over the lines. Adding a fourth statistic later means nesting one more both and extending tag's output tuple by one field, not adding a second loop.

\subsection*{10.4\quad Batch admission\hfill\textnormal{\textit{lab}}}
Admit a batch of form dicts with 7.3's mk\_user: all records, or the first rejection verbatim.

\noindent\texttt{admit(forms) -> Either [dict]}
\begin{lstlisting}[style=ex]
>>> admit([{'name': 'ada', 'age': '36'}, {'name': 'bob', 'age': '41'}])
('ok', [{'name': 'ada', 'age': 36}, {'name': 'bob', 'age': 41}])
>>> admit([{'name': 'ada', 'age': '36'}, {'name': '', 'age': '2'}])
('err', 'empty name')
\end{lstlisting}
\noindent\textbf{Solution.}
\begin{lstlisting}[style=code]
admit = o(sequenceE, partial(map, mk_user), list)
\end{lstlisting}
\noindent\textit{Note.} sequenceE, partial(map, mk\_user), and list are composed right to left with o: first every form in the batch runs through mk\_user's per-record validation from 7.3, producing a list of Ok/Err results; sequenceE then collapses that list into a single Either -- Ok of every validated record if all passed, or the first Err encountered if any one failed. The whole batch-intake pipeline fits on one line because each piece (mk\_user, sequenceE) was already built and verified independently in earlier chapters.

\clearpage
\part{Compare and Contrast: Imperative vs Functional}
\chapter*{About this part}
\addcontentsline{toc}{chapter}{About this part}
Ten small apps, each written twice over the same inputs: once imperative,
once with \texttt{haskell.py}, with a \texttt{\_\_main\_\_} block that
ASSERTS the two agree. Every listing below is verified to actually run and
agree before this book is built -- what you read is what was executed.

\begin{center}\small
\begin{tabular}{@{}p{0.13\textwidth}p{0.40\textwidth}p{0.38\textwidth}@{}}
\toprule
& \textbf{imperative} & \textbf{functional} \\
\midrule
state & mutable slots you maintain & threaded by the combinator \\
errors & exceptions / early return & values (NOTHING, Err) that compose \\
no value yet & hand-rolled sentinel per site & NOTHING, once, in the library \\
passes over data & one loop per concern & monoids fuse concerns into one pass \\
extension & edit inside the loop & add a stage / another \texttt{both} \\
\bottomrule
\end{tabular}
\end{center}
\chapter{calculator}
evaluate '2 * (3 + 4) - -5' both ways

\section*{Imperative}
\begin{lstlisting}[style=code]
def calc_imp(s):
    toks, i = [], 0                                   # 1. tokenize with an index
    while i < len(s):
        c = s[i]
        if c.isspace(): i += 1; continue
        if c in '+-*/()': toks.append(c); i += 1; continue
        if c.isdigit():
            j = i
            while j < len(s) and (s[j].isdigit() or s[j] == '.'): j += 1
            toks.append(float(s[i:j])); i = j; continue
        return ('err', f'bad char {c!r}')

    class P:                                          # 2. parse with a mutable cursor
        pos = 0
    def peek(): return toks[P.pos] if P.pos < len(toks) else None
    def unary():
        t = peek()
        if t == '-': P.pos += 1; return -unary()
        if t == '(':
            P.pos += 1; v = expr()
            if peek() != ')': raise ValueError('expected )')
            P.pos += 1; return v
        if isinstance(t, float): P.pos += 1; return t
        raise ValueError(f'unexpected {t!r}')
    def term():
        v = unary()
        while peek() in ('*', '/'):
            op = peek(); P.pos += 1; w = unary()
            v = v * w if op == '*' else v / w
        return v
    def expr():
        v = term()
        while peek() in ('+', '-'):
            op = peek(); P.pos += 1; w = term()
            v = v + w if op == '+' else v - w
        return v

    try:                                              # 3. errors ride exceptions
        v = expr()
        return ('ok', v) if P.pos == len(toks) else ('err', 'unconsumed input')
    except ValueError as e:
        return ('err', str(e))
\end{lstlisting}
\section*{Functional}
\begin{lstlisting}[style=code]
number = rx(r'-?\d+(\.\d+)?', float)
@doP
def _paren():
    yield lit('(')
    v = yield expr_fn
    yield lit(')')
    return v
unary_fn = lambda s: alt(number, _paren())(s)
term_fn  = chainl1(unary_fn, {'*': mul, '/': lambda a, b: a / b})
expr_fn  = chainl1(term_fn,  {'+': add, '-': sub})
calc_fn  = partial(runParser, expr_fn)
\end{lstlisting}
\section*{Discussion}

The imperative version's mutable cursor (\texttt{class P: pos = 0}) is a
common workaround for parsing in Python: state lives on an object because
a closure can't reassign an outer \texttt{int} without \texttt{nonlocal}
threaded through every helper. The functional version needs no cursor at
all -- \texttt{doP} threads the remaining input as an ordinary return
value, and \texttt{chainl1} handles left-associativity as a fold instead
of a hand-written \texttt{while} loop. Errors ride two different rails,
too: exceptions in the imperative version, an \texttt{Either} in the
functional one -- and only the \texttt{Either} can be inspected, logged,
or handed to another function without a \texttt{try}/\texttt{except} at
every call site.

\chapter{gradebook}
validate 'name,score' rows, then one-pass class stats

\section*{Imperative}
\begin{lstlisting}[style=code]
def admit_imp(rows):
    out = []                                          # accumulate + early return
    for row in rows:
        name, _, raw = row.partition(',')
        if not name:
            return ('err', f'empty name in {row!r}')
        if not raw.strip().isdigit():
            return ('err', f'bad score in {row!r}')
        score = int(raw)
        if score > 100:
            return ('err', f'score > 100 in {row!r}')
        out.append((name, score))
    n = tot = 0                                       # second loop for stats
    best = float('-inf')
    for _, s in out:
        n += 1; tot += s
        if s > best: best = s
    return ('ok', (out, (n, tot), best))
\end{lstlisting}
\section*{Object-Oriented}
\begin{lstlisting}[style=code]
class GradeBook:
    """The same job, encapsulated: a mutable object that grows by .add(row),
    raising on the first bad row instead of returning an Err tuple."""
    def __init__(self):
        self.rows = []

    def add(self, row):
        name, _, raw = row.partition(',')
        if not name:
            raise ValueError(f'empty name in {row!r}')
        if not raw.strip().isdigit():
            raise ValueError(f'bad score in {row!r}')
        score = int(raw)
        if score > 100:
            raise ValueError(f'score > 100 in {row!r}')
        self.rows.append((name, score))
        return self                                   # chainable: book.add(r1).add(r2)

    def stats(self):
        n = len(self.rows)
        tot = sum(s for _, s in self.rows)
        best = max((s for _, s in self.rows), default=float('-inf'))
        return (n, tot), best

def admit_oop(rows):
    book = GradeBook()
    try:
        for row in rows:
            book.add(row)
    except ValueError as e:
        return ('err', str(e))
    return ('ok', (book.rows, *book.stats()))
\end{lstlisting}
\section*{Functional}
\begin{lstlisting}[style=code]
@doE
def _row(row):
    name, _, raw = row.partition(',')
    yield (Ok(name) if name else Err(f'empty name in {row!r}'))
    s = yield note(f'bad score in {row!r}', int(raw) if raw.strip().isdigit() else NOTHING)
    yield (Ok(s) if s <= 100 else Err(f'score > 100 in {row!r}'))
    return Ok((name, s))

stats    = partial(foldMap, lambda p: ((1, p[1]), p[1]), both(both(Sum, Sum), MaxM))
admit_fn = lambda rows: bindE(traverseE(_row, rows),
                              lambda ps: Ok((ps, *stats(ps))))
\end{lstlisting}
\section*{Discussion}

Compare all three. The imperative version accumulates into a list and
returns early on a bad row, then needs a SECOND loop afterward just to
compute stats. The object-oriented version is the natural instinct for
validation-heavy code -- a class that grows by \texttt{.add()} and raises
on trouble -- but it still needs a \texttt{try}/\texttt{except} at the
call site, and its \texttt{.stats()} re-walks the same rows the loop
already built. The functional version fuses validation
(\texttt{traverseE}) and the one-pass fold (\texttt{foldMap} with
\texttt{both} nested twice) into two lines: one pass builds the valid
rows AND the stats at once, because \texttt{both} fuses monoids instead
of running loops back to back.

\chapter{layers}
layered config: defaults < file < cli; None is a real value

\section*{Imperative}
\begin{lstlisting}[style=code]
_UNSET = object()                                     # invent a private sentinel,
def effective_imp(layers, k):                         # because None must survive
    val = _UNSET
    for d in layers:
        if k in d:
            val = d[k]
    return NOTHING if val is _UNSET else val

def merged_imp(layers):
    out = {}
    for d in layers:
        for k, v in d.items():
            out[k] = v
    return out
\end{lstlisting}
\section*{Functional}
\begin{lstlisting}[style=code]
effective_fn = lambda layers, k: mconcat(Last, [maybe_get(d, k) for d in layers])
merged_fn    = lambda layers: foldl(lambda a, b: unionWith(lambda old, new: new, a, b),
                                    layers, {})
\end{lstlisting}
\section*{Discussion}

This is the one example where ``just use a config object'' would
recreate the exact bug the Maybe chapter exists to prevent: a plain class
doing \texttt{self.proxy = self.proxy or default} silently treats an
explicit \texttt{None} as ``no override'' -- backwards from what this
problem needs, since \texttt{FILECFG}'s \texttt{proxy: None} must WIN
over the default, not defer to it. The imperative fix (a private
\texttt{\_UNSET} sentinel) and the functional fix (\texttt{NOTHING}) are
the same idea; object-oriented encapsulation doesn't make the ambiguity
go away, it just relocates where you'd have to solve it.

\chapter{leaderboard}
clamp scores into [0, 100], rank by score desc, name asc

\section*{Imperative}
\begin{lstlisting}[style=code]
def rank_imp(rows):
    clamped = []
    for name, score in rows:
        if score < 0:
            score = 0
        elif score > 100:
            score = 100
        clamped.append((name, score))
    clamped.sort(key=lambda r: (-r[1], r[0]))          # score desc, name asc
    return clamped
\end{lstlisting}
\section*{Object-Oriented}
\begin{lstlisting}[style=code]
class Player:
    def __init__(self, name, score):
        self.name = name
        self.score = max(0, min(100, score))          # clamp at construction time

class Leaderboard:
    def __init__(self):
        self.players = []

    def add(self, name, score):
        self.players.append(Player(name, score))
        return self

    def ranked(self):
        ordered = sorted(self.players, key=lambda p: (-p.score, p.name))
        return [(p.name, p.score) for p in ordered]

def rank_oop(rows):
    board = Leaderboard()
    for name, score in rows:
        board.add(name, score)
    return board.ranked()
\end{lstlisting}
\section*{Functional}
\begin{lstlisting}[style=code]
clamp = lambda lo, hi: o(partial(max, lo), partial(min, hi))
clamp100 = clamp(0, 100)
name_of, score_of = itemgetter(0), itemgetter(1)
rank_fn = lambda rows: sortOn(
    lambda r: (-score_of(r), name_of(r)),
    [(name_of(r), clamp100(score_of(r))) for r in rows])
\end{lstlisting}
\section*{Discussion}

All three agree, so look at WHERE each one clamps. The imperative version
clamps inside the loop with an \texttt{if}/\texttt{elif}. The
object-oriented version clamps once, in \texttt{Player.\_\_init\_\_} --
so a \texttt{Player} object is a standing proof that its own score is
always valid, which is real encapsulation value, not just different
syntax for the same check. The functional version clamps by composing two
curried \texttt{max}/\texttt{min} calls into one reusable function
(\texttt{clamp}) -- no class needed to get the same ``validate once,
trust everywhere'' property, because \texttt{clamp100} IS the validated
boundary, not an object that happens to carry one.

\chapter{logtriage}
'LEVEL message' lines: error count, total, first error message

\section*{Imperative}
\begin{lstlisting}[style=code]
def triage_imp(lines):
    errors = 0                                        # three mutable slots
    total = 0
    first = None                                      # and a hand-rolled sentinel:
    seen_first = False                                # None alone can't mean "none yet"
    for ln in lines:                                  # if a message could be empty
        lvl, _, msg = ln.partition(' ')
        total += 1
        if lvl == 'ERROR':
            errors += 1
            if not seen_first:
                first, seen_first = msg, True
    return ((errors, total), first if seen_first else NOTHING)
\end{lstlisting}
\section*{Functional}
\begin{lstlisting}[style=code]
tag = lambda ln: (lambda lvl, _, msg:
                  ((lvl == 'ERROR', 1), msg if lvl == 'ERROR' else NOTHING)
                  )(*ln.partition(' '))
triage_fn = partial(foldMap, tag, both(both(Sum, Sum), First))
\end{lstlisting}
\section*{Discussion}

Two mutable slots plus a hand-rolled \texttt{seen\_first} flag in the
imperative version, because plain \texttt{None} can't distinguish ``no
error yet'' from ``the first error had an empty message''. The functional
version needs no flag at all: \texttt{First}'s identity is
\texttt{NOTHING}, not \texttt{None}, so the fold itself carries the
distinction. This is the same bug shape as \texttt{layers.py} --
\texttt{first}/\texttt{seen\_first} is a private sentinel invented from
scratch, precisely because Python's \texttt{None} is overloaded.

\chapter{orderbatch}
capstone: validate a batch of orders (qty > 0, price > 0), computing each line total; the whole batch fails at the first bad order, carrying which one and why

\section*{Imperative}
\begin{lstlisting}[style=code]
def total_imp(order):
    if order["qty"] <= 0:
        return ("err", f"{order['item']}: bad qty {order['qty']}")
    if order["price"] <= 0:
        return ("err", f"{order['item']}: bad price {order['price']}")
    return ("ok", round(order["qty"] * order["price"], 2))

def admit_imp(orders):
    out = []
    for o in orders:
        r = total_imp(o)
        if r[0] == "err":
            return r
        out.append(r[1])
    return ("ok", out)
\end{lstlisting}
\section*{Functional}
\begin{lstlisting}[style=code]
@doE
def total_fn(order):
    qty = yield (Ok(order["qty"]) if order["qty"] > 0
                 else Err(f"{order['item']}: bad qty {order['qty']}"))
    price = yield (Ok(order["price"]) if order["price"] > 0
                   else Err(f"{order['item']}: bad price {order['price']}"))
    return Ok(round(qty * price, 2))

admit_fn = lambda orders: sequenceE([total_fn(o) for o in orders])
\end{lstlisting}
\section*{Discussion}

Deliberately no object-oriented version here: an \texttt{Order} class
with a \texttt{.total()} method would just move \texttt{total\_imp}'s
body into a method, changing nothing about how a validation failure
PROPAGATES through a whole batch. That propagation -- stop at the first
bad order, carry which one and why -- is what \texttt{doE} and
\texttt{sequenceE} actually buy; encapsulating the per-order check
doesn't touch the batch-level control flow, which is the actual hard
part of this problem.

\chapter{orgchart}
resolve an employee's manager's manager's email, or NOTHING if any hop is missing. None-as-Nothing means a missing hop and a real missing-email must stay distinguishable

\section*{Imperative}
\begin{lstlisting}[style=code]
def skip_manager_email_imp(emp):
    mgr = MANAGER_OF.get(emp, NOTHING)
    if mgr is NOTHING:
        return NOTHING
    skip = MANAGER_OF.get(mgr, NOTHING)
    if skip is NOTHING:
        return NOTHING
    if skip not in EMAIL_OF:
        return NOTHING
    return EMAIL_OF[skip]
\end{lstlisting}
\section*{Functional}
\begin{lstlisting}[style=code]
@doM
def skip_manager_email_fn(emp):
    mgr = yield maybe_get(MANAGER_OF, emp)
    skip = yield maybe_get(MANAGER_OF, mgr)
    return (yield maybe_get(EMAIL_OF, skip))
\end{lstlisting}
\section*{Discussion}

Also deliberately no object-oriented version: wrapping
\texttt{MANAGER\_OF}/\texttt{EMAIL\_OF} in a class would just rename
\texttt{maybe\_get}/\texttt{doM} as methods, and the real challenge --
keeping ``employee not found'' distinguishable from ``employee found, no
email on file'' -- lives in the DATA MODEL, not the calling convention. A
hypothetical \texttt{.skip\_manager\_email()} written with plain
\texttt{dict.get(x)} returning \texttt{None} would silently reintroduce
the exact ambiguity \texttt{NOTHING} exists to prevent; you would just end
up reimplementing a sentinel inside the class anyway.

\chapter{payoff}
loan payoff schedule: balances after each fixed payment, and how many payments until paid off

\section*{Imperative}
\begin{lstlisting}[style=code]
def schedule_imp(balance, payment, rate):
    out, n = [balance], 0
    while balance > 0:
        balance = round(balance * (1 + rate) - payment, 2)
        if balance < 0:
            balance = 0.0
        out.append(balance)
        n += 1
    return out, n
\end{lstlisting}
\section*{Functional}
\begin{lstlisting}[style=code]
def _step(payment, rate):
    def step(balance):
        if balance <= 0:
            return NOTHING
        nxt = round(balance * (1 + rate) - payment, 2)
        nxt = max(nxt, 0.0)
        return (nxt, nxt)
    return step

def schedule_fn(balance, payment, rate):
    rest = unfoldr(_step(payment, rate), balance)
    return [balance] + rest, len(rest)
\end{lstlisting}
\section*{Discussion}

The imperative \texttt{while} loop and the functional \texttt{unfoldr}
compute the same recurrence, but \texttt{unfoldr}'s stopping condition is
DATA (\texttt{NOTHING}) returned from the step function, not a loop
condition checked from outside. That distinction matters most for
testing: \texttt{step(payment, rate)} is a pure function you can call
directly with one balance and inspect the answer, with no loop, no
mutable balance variable, and no need to run the whole schedule just to
check one step's arithmetic.

\chapter{sensoralert}
smooth a sensor stream with a 3-wide moving average, then report each consecutive run of over-threshold readings as (start, len)

\section*{Imperative}
\begin{lstlisting}[style=code]
def alerts_imp(xs, n, threshold):
    smoothed = []
    for i in range(len(xs) - n + 1):
        window = xs[i:i + n]
        smoothed.append(sum(window) / n)
    flags = [v > threshold for v in smoothed]
    runs, i = [], 0
    while i < len(flags):
        if flags[i]:
            start = i
            while i < len(flags) and flags[i]:
                i += 1
            runs.append((start, i - start))
        else:
            i += 1
    return runs
\end{lstlisting}
\section*{Functional}
\begin{lstlisting}[style=code]
def alerts_fn(xs, n, threshold):
    smoothed = [sum(w) / n for w in windows(n, xs)]
    flagged = [(i, v > threshold) for i, v in enum(smoothed)]
    groups = groupBy(lambda a, b: a[1] == b[1], flagged)
    return [(g[0][0], len(g)) for g in groups if g[0][1]]
\end{lstlisting}
\section*{Discussion}

\texttt{windows} and \texttt{groupBy} each replace a hand-rolled
index-walking loop with a name. The imperative version's inner
\texttt{while} (walking \texttt{flags} to find where each run ends) is
exactly what \texttt{groupBy} does generically -- once ``runs of the same
flag value'' is recognised as a \texttt{groupBy} problem, the run-finding
code disappears entirely, replaced by grouping and then filtering for the
runs that are \texttt{True}.

\chapter{wordfreq}
top-N word frequencies, ties alphabetical

\section*{Imperative}
\begin{lstlisting}[style=code]
def top_imp(n, text):
    counts = {}                                       # mutate a dict
    for w in text.lower().split():
        if w in counts:
            counts[w] += 1
        else:
            counts[w] = 1
    items = list(counts.items())                      # then sort, then slice
    items.sort(key=lambda kv: (-kv[1], kv[0]))
    return items[:n]
\end{lstlisting}
\section*{Functional}
\begin{lstlisting}[style=code]
top_fn = lambda n, text: pipe(
    text, str.lower, words,
    lambda ws: fromListWith(add, [(w, 1) for w in ws]),
    dict.items, partial(sortOn, lambda kv: (-kv[1], kv[0])),
    partial(take, n))
\end{lstlisting}
\section*{Discussion}

\texttt{pipe} reads top to bottom in the order the data actually flows --
lowercase, split, count, sort, take -- while the imperative version's
five separate steps (build counts, get items, sort in place, slice) are
the same five stages with no name tying them together. Neither version is
shorter by much; the functional one is easier to extend, because
inserting a new stage means inserting one more line in the \texttt{pipe}
call, not renumbering a sequence of loop variables.

\appendix
\part{Appendices}
\chapter{haskell.py, Complete}
\begin{lstlisting}[style=file]
"""haskell.py -- the comprehensive Haskell-flavoured toolbox for Python.

One importable vocabulary: the working parts of the Prelude, Data.List, Data.Maybe, Data.Either,
Data.Map, Data.Function, Control.Monad, monoids, and parser combinators, folded into sections and
mastered as one language.

Design rules
  * Where the stdlib already has it, the Haskell name WRAPS the stdlib (scanl is
    itertools.accumulate, memo is functools.cache) -- one vocabulary, zero reimplementation.
  * Failure is a value.  Maybe uses the true NOTHING sentinel, so None is LEGAL data everywhere;
    parsers fail with FAIL, and Either carries why.
  * do-notation via generators: `yield` is >>= (linear monads only).
  * Constant stack wherever Haskell would lean on TCO: loops, not recursion.
  * Python keywords and builtins keep prelude.py's underscore convention (not_, map_, filter_,
    zip_); `id` stays Python's -- use identity.
"""
from functools import reduce as _reduce, partial, cache as memo, cmp_to_key
from itertools import accumulate as _acc, islice as _islice, pairwise as _pairwise
from itertools import chain, count, cycle, repeat
from operator import add, sub, mul, not_, eq
from math import gcd, lcm, hypot, isqrt, prod as product
from collections import Counter, defaultdict, deque
from heapq import heappush, heappop, merge as _hmerge
import re as _re

# ============ 1. functions ============
identity = lambda x: x
const    = lambda x: lambda *_: x
flip     = lambda f: lambda a, b: f(b, a)
curry    = lambda f: lambda x: lambda y: f(x, y)
uncurry  = lambda f: lambda p: f(*p)

def compose(*fs):
    """Right-to-left composition, iterative (constant stack).
    >>> compose(str, abs)(-3)
    '3'
    >>> o(str, abs)(-3)
    '3'
    """
    def g(x):
        for f in reversed(fs):
            x = f(x)
        return x
    return g

o = compose                      # ML's operator: o(f, g) is f after g

def pipe(x, *fs):
    """Left-to-right application: pipe(x, f, g) == g(f(x)).
    >>> pipe(-3, abs, str)
    '3'
    """
    for f in fs:
        x = f(x)
    return x

def on(f, g):
    """Data.Function: combine two arguments via a key.
    >>> on(sub, len)('haskell', 'py')
    5
    """
    return lambda x, y: f(g(x), g(y))

# ============ 2. pairs ============
fst  = lambda p: p[0]
snd  = lambda p: p[1]
swap = o(tuple, reversed)

# ============ 3. sentinels ============
class _Sentinel:
    __slots__ = ("_n",)
    def __init__(self, n): self._n = n
    def __repr__(self):    return self._n
    def __bool__(self):    return False

NOTHING = _Sentinel("NOTHING")   # Maybe's Nothing; None is a legal Just here
FAIL    = _Sentinel("FAIL")      # parser failure; (None, rest) is a legal success
_MISS   = _Sentinel("_MISS")     # internal: absent optional argument

# ============ 4. arithmetic & logic ============
# add, sub, mul, not_ are operator's; gcd, lcm, hypot, isqrt, product are math's
div       = lambda a, b: a // b                       # floors, like Haskell div
mod       = lambda a, b: a % b                        # sign follows b
halve     = lambda n: n // 2
even      = lambda n: n % 2 == 0
odd       = lambda n: n % 2 == 1
otherwise = True

def quot(a, b):
    """Truncating division, Haskell quot; rem's sign follows a.
    >>> quotRem(-7, 2)
    (-3, -1)
    """
    return -(-a // b) if (a < 0) != (b < 0) else a // b
rem     = lambda a, b: a - b * quot(a, b)
quotRem = lambda a, b: (quot(a, b), rem(a, b))
signum  = lambda x: (x > 0) - (x < 0)
succ    = lambda x: chr(ord(x) + 1) if isinstance(x, str) else x + 1
pred    = lambda x: chr(ord(x) - 1) if isinstance(x, str) else x - 1

# ============ 5. folds, scans, unfolds ============
def foldl(f, xs, base=_MISS):
    """foldl; seeds with the first element if no base is given.
    >>> foldl(lambda a, b: a - b, [1, 2, 3], 10)
    4
    """
    return _reduce(f, xs) if base is _MISS else _reduce(f, xs, base)

def foldr(f, xs, base):
    """Right fold without recursion.
    >>> foldr(lambda x, acc: [x] + acc, [1, 2, 3], [])
    [1, 2, 3]
    """
    return _reduce(flip(f), reversed(list(xs)), base)

foldl1 = foldl
foldr1 = lambda f, xs: _reduce(flip(f), reversed(list(xs)))

def scanl(f, xs, base=_MISS):
    """All intermediate foldl states (itertools.accumulate).
    >>> scanl(add, [1, 2, 3], 0)
    [0, 1, 3, 6]
    """
    return list(_acc(xs, f) if base is _MISS else _acc(xs, f, initial=base))

scanl1 = lambda f, xs: scanl(f, xs)
scanr  = lambda f, xs, base: scanl(flip(f), reversed(list(xs)), base)[::-1]
scanr1 = lambda f, xs: scanl(flip(f), reversed(list(xs)))[::-1]

def unfoldr(f, seed):
    """Anamorphism: f(seed) -> (value, seed') or NOTHING to stop.
    >>> unfoldr(lambda n: NOTHING if n == 0 else (n, n - 1), 4)
    [4, 3, 2, 1]
    """
    out = []
    while (r := f(seed)) is not NOTHING:
        out.append(r[0]); seed = r[1]
    return out

def until(cond, f, x):
    """Loop f until cond holds.
    >>> until(lambda n: n > 100, lambda n: n * 2, 1)
    128
    """
    while not cond(x):
        x = f(x)
    return x

# brute-force enumerations (say the bound out loud):
replicateM   = lambda n, xs: foldl(lambda acc, _: [s + [x] for s in acc for x in xs], range(n), [[]])   # |xs|^n
subsequences = lambda xs: foldl(lambda acc, x: acc + [s + [x] for s in acc], list(xs), [[]])            # 2^n

# ============ 6. streams ============ (lazy; chain, count, cycle, repeat are itertools')
def iterate(f, x):
    """Infinite stream x, f(x), f(f(x)), ...
    >>> take(5, iterate(lambda n: n * 2, 1))
    [1, 2, 4, 8, 16]
    """
    while True:
        yield x
        x = f(x)

take = lambda n, xs: list(_islice(iter(xs), n))       # safe on infinite streams

# ============ 7. list basics ============
head        = lambda xs: xs[0]
tail        = lambda xs: xs[1:]
init        = lambda xs: xs[:-1]
last        = lambda xs: xs[-1]
drop        = lambda n, xs: list(xs)[max(n, 0):]
splitAt     = lambda n, xs: (xs[:n], xs[n:]) if n >= 0 else (xs[:0], xs)
replicate   = lambda n, x: [x] * n
elem        = lambda x, xs: x in xs
notElem     = lambda x, xs: x not in xs
null        = lambda xs: len(xs) == 0
nub         = o(list, dict.fromkeys)                  # unique, first seen wins
enum        = lambda xs, start=0: list(enumerate(xs, start))
pairwise    = lambda xs: list(_pairwise(xs))
zip_        = lambda a, b: list(zip(a, b))
zip3        = lambda a, b, c: list(zip(a, b, c))
unzip       = lambda ps: tuple(map(list, zip(*ps))) if ps else ([], [])
isPrefixOf  = lambda p, xs: list(xs[:len(p)]) == list(p)
isSuffixOf  = lambda p, xs: list(xs[len(xs) - len(p):]) == list(p)

def lookup(k, pairs):
    """Assoc-list lookup into Maybe (a stored None is data).
    >>> lookup('b', [('a', 1), ('b', None)]) is None
    True
    >>> lookup('c', [('a', 1)]) is NOTHING
    True
    """
    return next((v for kk, v in pairs if kk == k), NOTHING)

def stripPrefix(p, xs):
    """The rest after prefix p, or NOTHING.
    >>> stripPrefix('foo', 'foobar')
    'bar'
    >>> stripPrefix('x', 'foobar') is NOTHING
    True
    """
    return xs[len(p):] if isPrefixOf(p, xs) else NOTHING

# ============ 8. higher-order lists ============
map_      = lambda f, xs: [f(x) for x in xs]
filter_   = lambda cond, xs: [x for x in xs if cond(x)]
concat    = lambda xss: [x for xs in xss for x in xs]
concatMap = lambda f, xs: [y for x in xs for y in f(x)]
cross     = lambda a, b: [(x, y) for x in a for y in b]
starmap   = lambda f, pairs: [f(*p) for p in pairs]
zipWith   = lambda f, a, b: [f(x, y) for x, y in zip(a, b)]
zipWith3  = lambda f, a, b, c: [f(x, y, z) for x, y, z in zip(a, b, c)]

def find(cond, xs):
    """First match, lazily, into Maybe.
    >>> find(even, [1, 3, 4, 5])
    4
    >>> find(even, [1, 3]) is NOTHING
    True
    """
    return next((x for x in xs if cond(x)), NOTHING)

def partition(cond, xs):
    """(keepers, rest) in ONE pass, cond called once per element.
    >>> partition(even, [1, 2, 3, 4])
    ([2, 4], [1, 3])
    """
    yes, no = [], []
    for x in xs:
        (yes if cond(x) else no).append(x)
    return (yes, no)

# ============ 9. slicing & spans ============
def span(cond, xs):
    """Split at first failure, one pass, one cond call per element.
    >>> span(lambda x: x < 3, [1, 2, 3, 1])
    ([1, 2], [3, 1])
    """
    xs = list(xs)
    i = next((i for i, x in enumerate(xs) if not cond(x)), len(xs))
    return (xs[:i], xs[i:])

break_ = lambda cond, xs: span(o(not_, cond), xs)

def takewhile(cond, xs):                    # generator
    for x in xs:
        if not cond(x):
            return
        yield x

def dropwhile(cond, xs):                    # generator
    it = iter(xs)
    for x in it:
        if not cond(x):
            yield x
            break
    yield from it

def takeuntil(cond, xs):                    # takewhile that INCLUDES the first hit
    for x in xs:
        yield x
        if cond(x):
            return

takeWhile = lambda cond, xs: list(takewhile(cond, xs))
dropWhile = lambda cond, xs: list(dropwhile(cond, xs))
inits     = lambda xs: [xs[:i] for i in range(len(xs) + 1)]
tails     = lambda xs: [xs[i:] for i in range(len(xs) + 1)]

def chunksOf(n, xs):
    """
    >>> chunksOf(2, [1, 2, 3, 4, 5])
    [[1, 2], [3, 4], [5]]
    """
    xs = list(xs)
    return [xs[i:i + n] for i in range(0, len(xs), n)]

def windows(n, xs):
    """Sliding windows of width n.
    >>> windows(3, [1, 2, 3, 4])
    [[1, 2, 3], [2, 3, 4]]
    """
    xs = list(xs)
    return [xs[i:i + n] for i in range(len(xs) - n + 1)]

def groupBy(eq, xs):
    """Haskell semantics: runs anchored on each run's FIRST element.
    >>> groupBy(lambda a, b: a == b, "aabba")
    [['a', 'a'], ['b', 'b'], ['a']]
    """
    out = []
    for x in xs:
        if out and eq(out[-1][0], x):
            out[-1].append(x)
        else:
            out.append([x])
    return out

group = partial(groupBy, eq)

def transpose(rows):
    """Ragged-safe like Data.List.
    >>> transpose([[1, 2, 3], [4, 5]])
    [[1, 4], [2, 5], [3]]
    """
    rows = [list(r) for r in rows]
    out, i = [], 0
    while any(i < len(r) for r in rows):
        out.append([r[i] for r in rows if i < len(r)])
        i += 1
    return out

# ============ 10. sorting & searching ============
sortOn = lambda f, xs: sorted(xs, key=f)
sortBy = lambda cmp, xs: sorted(xs, key=cmp_to_key(cmp))
maxOn  = lambda f, xs: max(xs, key=f)
minOn  = lambda f, xs: min(xs, key=f)
merge  = lambda a, b: list(_hmerge(a, b))             # stable merge of sorted lists

# ============ 11. strings ============
lines   = str.splitlines
unlines = lambda ls: "".join(l + "\n" for l in ls)
words   = str.split
unwords = " ".join

# ============ 12. containers ============ (Counter, defaultdict, deque, heappush/heappop are stdlib's)
def fromListWith(f, pairs):
    """Data.Map: THE dict-building fold; f(new, old).
    >>> fromListWith(add, [('a', 1), ('b', 2), ('a', 3)]) == {'a': 4, 'b': 2}
    True
    """
    d = {}
    for k, v in pairs:
        d[k] = f(v, d[k]) if k in d else v
    return d

def insertWith(f, k, v, d):
    """Persistent insert (fresh dict); f(new, old).
    >>> insertWith(add, 'a', 5, {'a': 1}) == {'a': 6}
    True
    """
    return {**d, k: f(v, d[k]) if k in d else v}

def unionWith(f, a, b):
    """Data.Map: merge two dicts, combining shared keys with f.
    >>> unionWith(add, {'a': 1, 'b': 2}, {'b': 3, 'c': 4}) == {'a': 1, 'b': 5, 'c': 4}
    True
    """
    return {k: f(a[k], b[k]) if k in a and k in b else (a[k] if k in a else b[k])
            for k in a.keys() | b.keys()}

def getpath(t, ks, default=NOTHING):
    """Read along a key path without autovivifying; Maybe out.
    >>> getpath({'a': {'b': None}}, ['a', 'b']) is None
    True
    >>> getpath({'a': {}}, ['a', 'b']) is NOTHING
    True
    """
    for k in ks:
        if not isinstance(t, dict) or k not in t:
            return default
        t = t[k]
    return t

# multiset ops on count-dicts (use .get -- indexing a Counter autovivifies):
bag_diff  = lambda a, b: {k: a[k] - b.get(k, 0) for k in a if a[k] > b.get(k, 0)}
bag_inter = lambda a, b: {k: v for k in a.keys() & b.keys() if (v := min(a[k], b[k]))}
bag_union = lambda a, b: {k: max(a.get(k, 0), b.get(k, 0)) for k in a.keys() | b.keys()}
bag_sub   = lambda a, b: all(b.get(k, 0) >= n for k, n in a.items())

# ============ 13. Maybe ============ (NOTHING sentinel; None is a legal value)
def bind(x, f):
    """>>= : short-circuit on NOTHING.
    >>> bind(3, lambda v: v + 1), bind(NOTHING, lambda v: v + 1)
    (4, NOTHING)
    """
    return NOTHING if x is NOTHING else f(x)

isJust    = lambda x: x is not NOTHING
isNothing = lambda x: x is NOTHING
fromMaybe = lambda default, x: default if x is NOTHING else x

listToMaybe = lambda xs: xs[0] if xs else NOTHING
maybeToList = lambda x: [] if x is NOTHING else [x]

def catMaybes(xs):
    """
    >>> catMaybes([1, NOTHING, None, 2])
    [1, None, 2]
    """
    return [x for x in xs if x is not NOTHING]

def mapMaybe(f, xs):
    """Map into Maybe and drop the NOTHINGs, one pass.
    >>> mapMaybe(lambda s: int(s) if s.isdigit() else NOTHING, ['3', 'x', '7'])
    [3, 7]
    """
    return [y for y in map(f, xs) if y is not NOTHING]

def sequenceM(xs):
    """All-or-nothing.
    >>> sequenceM([1, 2]), sequenceM([1, NOTHING])
    ([1, 2], NOTHING)
    """
    out = []
    for x in xs:
        if x is NOTHING:
            return NOTHING
        out.append(x)
    return out

traverseM = lambda f, xs: sequenceM(map_(f, xs))

def maybe_get(d, k):
    """dict lookup into Maybe (None stays distinguishable from missing).
    >>> maybe_get({'a': None}, 'a'), maybe_get({}, 'a')
    (None, NOTHING)
    """
    return d.get(k, NOTHING)

# ============ 14. Either ============ (tagged tuples; match-friendly)
Ok  = lambda v: ("ok", v)
Err = lambda e: ("err", e)

def bindE(x, f):
    """
    >>> bindE(Ok(3), lambda v: Ok(v + 1)), bindE(Err("no"), lambda v: Ok(v))
    (('ok', 4), ('err', 'no'))
    """
    return x if x[0] == "err" else f(x[1])

def sequenceE(xs):
    """First Err wins; else Ok of all values.
    >>> sequenceE([Ok(1), Err("a"), Err("b")])
    ('err', 'a')
    """
    out = []
    for x in xs:
        if x[0] == "err":
            return x
        out.append(x[1])
    return Ok(out)

def traverseE(f, xs):
    """Map into Either, then sequence.
    >>> traverseE(lambda n: Ok(n * 2) if n else Err('zero'), [1, 2])
    ('ok', [2, 4])
    """
    return sequenceE(map_(f, xs))

def partitionEithers(xs):
    """
    >>> partitionEithers([Ok(1), Err('a'), Ok(2)])
    ([1, 2], ['a'])
    """
    oks, errs = [], []
    for tag, v in xs:
        (oks if tag == "ok" else errs).append(v)
    return (oks, errs)

def note(msg, x):
    """Maybe -> Either: attach WHY to a NOTHING.
    >>> note("missing key", NOTHING)
    ('err', 'missing key')
    """
    return Err(msg) if x is NOTHING else Ok(x)

# ============ 15. do-notation ============ (yield is >>=; linear monads only)
def do(binder, is_fail):
    """Build a do-block decorator for any short-circuiting monad."""
    def deco(gen_fn):
        def run(*a, **kw):
            g = gen_fn(*a, **kw)
            def step(val):
                try:
                    m = g.send(val)
                except StopIteration as e:
                    return e.value
                return binder(m, step)
            return step(None)
        return run
    return deco

doM = do(bind,  lambda x: x is NOTHING)     # generator returns a plain value
doE = do(bindE, lambda x: x[0] == "err")    # generator returns Ok(...)/Err(...)

# ============ 16. parser combinators ============
# Parser a = str -> (value, rest) | FAIL.   Values may be None (json null!).
def doP(gen_fn):
    """do-notation that THREADS the input: yield a parser, receive its value.
    Fresh generator per parse -- re-runnable, backtrack-safe under alt."""
    def make(*a, **kw):
        def parser(s):
            g, val = gen_fn(*a, **kw), None
            try:
                while True:
                    r = g.send(val)(s)
                    if r is FAIL:
                        return FAIL
                    val, s = r
            except StopIteration as e:
                return (e.value, s)
        return parser
    return make

def alt(*ps):
    """<|> : first success wins."""
    return lambda s: next((r for r in (p(s) for p in ps) if r is not FAIL), FAIL)

def many(p):
    """Zero or more, loop not recursion; stalls guard against empty matches."""
    def parser(s):
        out = []
        while (r := p(s)) is not FAIL and r[1] != s:
            out.append(r[0]); s = r[1]
        return (out, s)
    return parser

def many1(p):
    def parser(s):
        r = p(s)
        if r is FAIL:
            return FAIL
        rest = many(p)(r[1])
        return ([r[0]] + rest[0], rest[1])
    return parser

def sepBy(p, sep):
    def pair(s):
        r = sep(s)
        return FAIL if r is FAIL else p(r[1])
    def parser(s):
        r = p(s)
        if r is FAIL:
            return ([], s)
        rest = many(pair)(r[1])
        return ([r[0]] + rest[0], rest[1])
    return parser

def lit(c):
    """One expected character, whitespace-blind."""
    def parser(s):
        s = s.lstrip()
        return (c, s[1:]) if s[:1] == c else FAIL
    return parser

def rx(pat, conv=identity):
    """Regex token, whitespace-blind, converted."""
    r = _re.compile(pat)
    def parser(s):
        s = s.lstrip()
        m = r.match(s)
        return (conv(m.group()), s[m.end():]) if m else FAIL
    return parser

def chainl1(p, ops):
    """p (op p)* folded LEFT; ops maps operator char -> binary function."""
    opp = alt(*[lit(c) for c in ops])
    def parser(s):
        vr = p(s)
        while vr is not FAIL:
            op = opp(vr[1])
            if op is FAIL:
                return vr
            w = p(op[1])
            if w is FAIL:
                return vr
            vr = (ops[op[0]](vr[0], w[0]), w[1])
        return FAIL
    return parser

def runParser(p, s):
    """Parser -> Either: Ok(value) or Err showing where it stopped.
    >>> runParser(rx(r'\\d+', int), "42")
    ('ok', 42)
    >>> runParser(rx(r'\\d+', int), "42x")
    ('err', "unconsumed input: 'x'")
    """
    r = p(s)
    if r is FAIL:
        return Err(f"no parse: {s!r}")
    if r[1].strip():
        return Err(f"unconsumed input: {r[1].strip()!r}")
    return Ok(r[0])

# ============ 17. monoids ============
Monoid  = lambda empty, op: (empty, op)
mconcat = lambda m, xs: foldl(m[1], xs, m[0])
foldMap = lambda f, m, xs: mconcat(m, [f(x) for x in xs])
both    = lambda m1, m2: Monoid((m1[0], m2[0]),
              lambda a, b: (m1[1](a[0], b[0]), m2[1](a[1], b[1])))
Sum     = Monoid(0, add)
Product = Monoid(1, mul)
All     = Monoid(True,  lambda a, b: a and b)
Any     = Monoid(False, lambda a, b: a or b)
MaxM    = Monoid(float("-inf"), max)
MinM    = Monoid(float("inf"),  min)
ListM   = Monoid([], add)
First   = Monoid(NOTHING, lambda a, b: b if a is NOTHING else a)
Last    = Monoid(NOTHING, lambda a, b: a if b is NOTHING else b)

if __name__ == "__main__":
    import doctest
    fails, total = doctest.testmod()
    print(f"{total - fails}/{total} doctests passed" if not fails
          else f"{fails} FAILED")

\end{lstlisting}
\backmatter
\end{document}
